\pdfoutput=1
\documentclass[11pt,reqno]{amsart}
\usepackage[top=1in,bottom=1in,left=1.2in,right=1.2in]{geometry}
\usepackage[T1]{fontenc}
\usepackage[utf8]{inputenc}
\usepackage{amsmath,amsthm,amssymb,amsfonts}
\usepackage{mathtools}
\usepackage{hyperref}
\hypersetup{hidelinks}
\usepackage{enumitem}
\usepackage{lmodern}
\usepackage{xcolor}
\usepackage[protrusion=true,expansion=false]{microtype} % - Theorem environments -
\setlength{\emergencystretch}{2em} % avoid overfull lines
\usepackage{caption}
% Revision markup removed
\newtheorem{theorem}{Theorem}[section]
\newtheorem{proposition}[theorem]{Proposition}
\newtheorem{lemma}[theorem]{Lemma}
\newtheorem{corollary}[theorem]{Corollary}
\theoremstyle{definition}
\newtheorem{definition}[theorem]{Definition}
\theoremstyle{remark}
\newtheorem{remark}[theorem]{Remark} % - Notation -
\newcommand{\R}{\mathbb{R}}
\newcommand{\Rp}{\mathbb{R}_{>0}}
\newcommand{\1}{\mathbf{1}}
\newcommand{\norm}[1]{\left\lVert #1 \right\rVert}
\newcommand{\avg}[1]{\overline{#1}}
\newcommand{\zero}{\mathrm{zero}}
\newcommand{\nonzero}{\mathrm{nonzero}}
\newcommand{\inconclusive}{\mathrm{inconclusive}}
\title[Coercive Projection Theorem]{The Coercive Projection Theorem for\\
Canonical Reciprocal Costs}
\author{Jonathan Washburn}
\address{Recognition Physics Research Institute, Austin, TX USA}
\email{jon@recognitionphysics.org}

\author{Amir Rahnamai Barghi}
\address{Recognition Physics Research Institute, Austin, TX USA}
\email{arahnamab@gmail.com}
\thanks{* Correspondence: \texttt{arahnamab@gmail.com}.}

\date{} % Added for measurement-matrix section
\DeclareMathOperator{\rank}{rank} % use operator form; avoid macro conflicts
\newcommand{\Mean}{\operatorname{Mean}}
% - Notation macros (consistency) -
\providecommand{\C}{\mathbb{C}}
\providecommand{\N}{\mathbb{N}}
\providecommand{\Z}{\mathbb{Z}}
\newcommand{\eps}{\varepsilon}

% Hypothesis tags (reader-friendly)
\newcommand{\Hrank}{\textbf{(H-rank)}}
\newcommand{\Hprony}{\textbf{(H-prony)}}
\newcommand{\Hrec}{\textbf{(H-rec)}}
\newcommand{\Hneu}{\textbf{(H-neu)}}
\newcommand{\Hsep}{\textbf{(H-sep)}}
\begin{document}
% Abstract must precede \maketitle in amsart.
\begin{abstract}
We develop a finite-data framework for certifying \emph{zero-defect} (neutral) configurations of positive vectors under the canonical separable reciprocal cost.
We show that this scalar cost is characterized among nonconstant continuous costs by the Recognition Composition Law together with a local quadratic calibration at balance; in particular, reciprocity symmetry and the normalization at the neutral point follow from the composition law.
Under a conservation constraint and short window observations of a rational (finite--state) signal class, we construct a canonical decision procedure that is \emph{locally maximal on the identifiability locus} among all sound procedures: any sound rule that resolves a datum must agree with the canonical output, and cannot resolve strictly more cases.

The method is organized as $\Phi^\ast=A\circ B\circ P$: the projection/coercivity core is forced by the canonical-cost axioms, while the aggregation/reconstruction step is specified on a nondegenerate identifiability locus (e.g.\ a Hankel invertibility condition).

\end{abstract}
\maketitle
\par\medskip
\noindent\textbf{MSC 2020:} Primary 90C25; Secondary 47H05, 94A12.\par
\noindent\textbf{Keywords:} canonical reciprocal cost; convex analysis; proximal maps; coercivity; finite-window aggregation; inverse problems; certification.\par
\medskip

%\maketitle moved above abstract for MDPI-style placement.
\section{Introduction}

We introduce a finite-window certification problem in which observed aggregates are tested for compatibility with a neutral \emph{zero-defect} class. After fixing notation and the ratio-induced cost framework, we summarize the main results and explain how the later sections build the proof pipeline. For additional background on the broader Recognition Science program and its geometric framing, see \cite{washburnnetal2026RG}.

\paragraph{Relation to Hankel/Prony reconstruction.} Although our setting is organized around aggregated window data and a canonical convex-analytic decision core, the nondegenerate identifiability locus is expressed via finite Hankel full-rank/invertibility conditions. This connects to the classical finite-rank/Hankel-structure viewpoint (e.g.\ Kronecker-type characterizations and related operator theory) and to Prony-style exponential fitting, where solving structured linear systems built from short samples is central; see, for example, \cite{peller2003hankel,osborne1995prony}.

\subsection{Connection to classical projection and proximal tools}
The certification template is \emph{forced} by our axioms, but its shape echoes classical projection/proximal
ideas from convex analysis. We refer to the main theorem as the \emph{Coercive Projection Theorem} (CPT), and we
write its canonical finite-data certifier as $\Phi^\ast$. First, the CPT certifier $\Phi^\ast$ begins with a projection-like update built from the
canonical reciprocal cost together with a coercive regularization; at a structural level this plays the role of
a proximity map that turns a penalty into a stable ``best-approximation'' operator (cf. Moreau's proximity
operator and envelope \cite{moreau1965}). Second, because our state variables are inherently multiplicative
(ratios/scales), the natural geometry is non-Euclidean: Bregman divergences \cite{bregman1967} describe projection in coordinates adapted to log/ratio structure, matching the way our axioms select log coordinates as
the intrinsic ``neutral'' parameterization. Third, the coercive resolvent-type step that converts defect bounds
into uniqueness and perturbation stability is in the spirit of Rockafellar's proximal-point framework for
monotone operators \cite{rockafellar1976}.

The difference is that the classical theory typically offers a \emph{family} of legitimate proximal/projection
schemes (choice of divergence, step size, regularizer). In this paper we separate two levels explicitly: Lemma~\ref{lem:forced-factor} gives the per-state order-theoretic representation, and Theorem~\ref{thm:forced-factor-unique} gives the global uniqueness/forcedness statement under explicit cross-state rigidity hypotheses. The subsequent aggregation/reconstruction stage $A$ is a standard inverse-problem ingredient and is only invoked on a nondegenerate (locally invertible) measurement locus; it is an additional hypothesis beyond the cost axioms. This separation is what makes the resulting certifier both canonical (in its cost-driven components) and reproducible.

\subsection{The problem}
Given finite observational access to a configuration $x\in(\Rp)^n$, decide whether it has
\emph{zero defect}. In the present setting, ``zero defect'' means membership in the structured
set
\[
S:=\{x\in(\Rp)^n:\ J_n(x)=0\}.
\]
\noindent Here \(J_n:(\Rp)^n\to[0,\infty)\) denotes the separable extension of \(J\),
\begin{equation}\label{eq:Jn}
J_n(x):=\sum_{i=1}^n J(x_i).
\end{equation}
More generally, ratio-induced recognition costs take the form
\begin{equation}\label{eq:cso}
c(s,o):=J\!\left(\frac{\iota_S(s)}{\iota_O(o)}\right),
\end{equation}
where $\iota_S,\iota_O$ are positive scale maps and $J:(0,\infty)\to[0,\infty)$ is a penalty.

Under inversion symmetry, strict convexity/coercivity, normalization at $1$, and a multiplicative
d'Alembert identity, one can classify admissible penalties and recover (up to a scale parameter) the same
canonical form $J(x)=\cosh(a\log x)-1$; see \cite{washburnbarghi2026RCC}. In this paper we take the
normalized choice \eqref{eq:Jdef} as given and focus on the finite-data certification problem. Under the canonical cost, $S$ collapses to the single configuration $x\equiv 1$ (all components
equal to $1$). The question is therefore the most basic decision task: determine from finite data
whether $x$ is exactly the neutral configuration.
\subsection{Main goal and results}

Our goal is to characterize the \emph{locally optimal} and \emph{structurally forced} certification procedure for the basic decision question ``is the configuration equal to the neutral state $x\equiv 1$?'' on the rational/finite-state class. Under three explicit hypotheses-(i) the canonical reciprocal cost forced by the composition law, (ii) a conservation constraint, and (iii) finite local resolution with window length $W\ge 1$ (e.g. $W=8$)-we show that the canonical procedure is locally maximal (in resolve-set domination) on the identifiability locus, and that any sound procedure that resolves there must agree with its output. Moreover, the operational pipeline uses three universal components: a projection step $P$, a coercivity step $B$, and an aggregation step $A$. The \emph{structure} of this decomposition and the sharp coercivity constant in $B$ are consequences of the canonical-cost axioms, while the reconstruction/aggregation step $A$ requires an \emph{additional} nondegeneracy (local invertibility) hypothesis for the window map (e.g. Hankel/Vandermonde nonsingularity as in Theorem~\ref{thm:window-ident-global}). Our maximality statements (e.g. Theorem~\ref{thm:domination}) are correspondingly local to that locus and do not claim global parameter recovery from window sums.

\subsection{Rigidity and local optimality of the canonical pipeline}
The term \emph{method} suggests optional design choices. In our setting we prove a rigidity statement: among all sound certification rules (Definition~\ref{def:sound}), the canonical pipeline $\Phi^\ast=A\circ B\circ P$ is locally maximal on the identifiability locus (Theorem~\ref{thm:domination}).

\paragraph{What is proved.}
Rather than packaging the contribution as a list of broad meta-properties, we summarize the concrete statements established in this paper:
\begin{enumerate}[label=\textup{(\roman*)},leftmargin=2.2em]
\item \textbf{Quantitative sound certification under noise.} Theorems~\ref{thm:eps-tolerant} and~\ref{thm:eps-cert} give explicit $\varepsilon$--stability bounds for the realized costs and the associated meaning sets, including the sharp two-sided tolerance (membership in $\Mean_{2\varepsilon}$ rather than $\Mean_{\varepsilon}$).
\item \textbf{Local domination among sound rules.} Theorem~\ref{thm:domination} proves that on the identifiability locus the canonical pipeline $\Phi^\ast=A\circ B\circ P$ is locally maximal among all sound certification rules (Definition~\ref{def:sound}), i.e.\ no sound rule can resolve strictly more instances in a neighborhood.
\item \textbf{Rigidity/uniqueness consequences of the axioms.} Theorem~\ref{thm:RCL-unique} characterizes the canonical reciprocal scalar cost from the Recognition Composition Law together with the local calibration at balance, and Theorem~\ref{thm:forced-factor-unique} records the induced factorization/uniqueness behavior for certificates under explicit rigidity hypotheses.
\item \textbf{Uniqueness among locally optimal rules.} Corollary~\ref{cor:unique-local-optimal} shows that on the identifiability neighborhood $\mathcal{U}_{d,W}$ the canonical rule is not only locally maximal but also unique among procedures achieving local maximality: any sound rule that is locally optimal on $\mathcal{U}_{d,W}$ must agree with $\Phi^\ast$ wherever it resolves.
\end{enumerate}



Lemma~\ref{lem:forced-factor} shows the local ordinal representation, while Theorem~\ref{thm:forced-factor-unique} proves uniqueness of the realized-cost reparametrization and of the state-free profile under explicit rigidity assumptions; the aggregation/reconstruction step $A$ is then pinned down on any nondegenerate (locally invertible) measurement locus, but is \emph{not} implied by the cost axioms alone. For that reason we treat
the procedure as a theorem-the \emph{Coercive Projection Theorem} (CPT)-rather than as an
algorithmic ``recipe'' one could freely modify.

We study a compatibility certificate for the neutral class: under the stated axioms, we give sufficient conditions ensuring that the observed finite-window aggregates admit a model-consistent realization in the neutral \emph{zero-defect} class. Criteria for ruling out neutrality, or for statistical hypothesis testing, are not treated here.

\begin{remark}\label{rem:kernel-dict}
For convenience we record a short dictionary between the objects defined in this paper and standard cost-based selection terminology. (We avoid the word ``kernel'' here, since in other contexts it denotes a nullspace or vanishing set.) The \emph{base cost} is our ratio-induced map $c:S\times O\to[0,\infty)$ given by $c(s,o):=J(\iota_S(s)/\iota_O(o))$. The associated \emph{meaning set} (argmin selector) is $\Mean(s):=\arg\min_{o\in O} c(s,o)$. The corresponding \emph{defect} (or \emph{gap}) is
\[
\Delta(s,o):=c(s,o)-\inf_{o'\in O}c(s,o'),
\]
so that $\Delta(s,o)=0$ iff $o\in\Mean(s)$. Finally, the CPT \emph{certificate} $\mathcal{C}$ is a scalar test statistic built from the $W$-block window measurements; in Section~\ref{sec:applications} it is used to separate the null (neutral) configuration from the nonzero class under the stated coercivity and identifiability hypotheses. This remark is purely notational and introduces no additional assumptions.
\end{remark}

\section{Model, axioms, and finite-window measurements}\label{sec:axioms}

This section sets up the mathematical model. We define the state and observation spaces, scale maps, the canonical reciprocal cost and induced defect, and the meaning/selection map. We then define the finite-window aggregation operator and list the axioms used in all subsequent theorems.

\medskip
\noindent\textbf{Big-picture guide.} Section~2 fixes the canonical cost and the neutral/defect geometry in log-coordinates (the forced steps $P$ and $B$). It then defines the finite-window measurement operator $\mathcal W$ and records the identifiability locus needed for the reconstruction step $A$. Section~3 uses these ingredients to define the finite-data procedure and prove soundness and domination (maximality) on the stated nondegenerate locus.
\medskip

We state the axioms and fix notation for later use. Proofs are provided where the results are stated. \begin{definition}\label{def:system}
A \emph{recognition cost system} is a quadruple $(\Rp,J,\sigma,W)$ where:
\begin{enumerate}[label=\textup{(C\arabic*)},leftmargin=2.2em]
\item \textbf{Recognition Composition Law (primitive axiom for $J$).}
We assume a scalar cost $J:(0,\infty)\to\mathbb{R}$ is non-constant, twice differentiable in a neighborhood of $1$, and satisfies the Recognition Composition Law (RCL):
This $J$ is the penalty appearing in the ratio-induced recognition cost $c(s,o)$ in \eqref{eq:cso}, and we write $J_n$ for its separable extension \eqref{eq:Jn}.
\begin{equation}\label{eq:RCL}
J(xy)+J(x/y)=2J(x)+2J(y)+2J(x)J(y),\qquad x>0,\ y>0,
\end{equation}
together with the curvature normalization $J''\!(1)=1$.
We use the corresponding canonical reciprocal cost,
\begin{equation}\label{eq:Jdef}
J(x):=\frac12\bigl(x+x^{-1}\bigr)-1,\qquad x>0.
\end{equation}
\item \textbf{Conservation.} For $x\in(\Rp)^n$, define
$\sigma(x):=\sum_{i=1}^n \log x_i$. Admissible configurations satisfy $\sigma(x)=0$.
\item \textbf{Finite resolution.} The observation window has fixed integer length $W\ge 1$ (in examples we take $W=8$):
only $W$ consecutive evaluations per cycle of the underlying finite-state representation
are available.
\end{enumerate}
\end{definition}

\noindent For later use we write $J_n(x):=\sum_{i=1}^n J(x_i)$ for $x\in(\Rp)^n$.

\begin{remark}
(i) The Recognition Composition Law together with non-constancy forces the balance normalization $J(1)=0$ and reciprocity $J(x)=J(1/x)$. Under the stated $C^2$ regularity near $1$ and the local quadratic calibration in Axiom~2 (equivalently $J''(1)=1$; see Appendix), the cost is uniquely determined; equivalently, it is given by \eqref{eq:Jdef}.

(ii) Section~\ref{subsec:domination} will make precise the ``rational/finite-state class'' and the
window data model used by the aggregation step.
\end{remark}

\subsection{Roadmap of the proof}\label{subsec:roadmap}
We will (i) identify the canonical cost and its log-coordinate form,
(ii) define the three operations $P,B,A$, and (iii) prove the local representation lemma (Lemma~\ref{lem:forced-factor}) together with the strengthened global uniqueness theorem (Theorem~\ref{thm:forced-factor-unique}) for the forced-factorization layer under explicit rigidity hypotheses.
Complete proofs of the window-data identifiability and domination steps are included in Sections~\ref{sec:aggregation} and \ref{subsec:domination}. % - BEGIN: added sections (v2) -

\subsection{Canonical cost and defect geometry}\label{sec:geometry} This subsection collects basic properties of the canonical reciprocal cost $J$ and the associated defect in log-coordinates, which will later control the aggregation and certification steps.

\subsubsection{The canonical reciprocal cost}
We recall that $J$ is fixed by \eqref{eq:Jdef}. The following logarithmic-coordinate reformulation is used repeatedly.
\begin{lemma}\label{lem:logform}
Let $\varphi:\R\to\R_{\geq 0}$ be defined by $\varphi(t):=\cosh(t)-1$.
Then for all $t\in\R$,
\[
J(e^t)=\varphi(t)=\cosh(t)-1.
\]
In particular, $J(x)=0$ if and only if $x=1$.
\end{lemma}

\begin{proof}
Using $\cosh(t)=\tfrac12(e^t+e^{-t})$ we compute
$J(e^t)=\tfrac12(e^t+e^{-t})-1=\cosh(t)-1$.
Since $\cosh(t)\geq 1$ with equality only at $t=0$, we have $J(e^t)=0\iff t=0\iff e^t=1$.
\end{proof}
\subsubsection{Neutrality and projection in log-coordinates}
We work in log-coordinates $y=\log x\in\R^n$ for $x\in(\Rp)^n$.
\begin{definition}\label{def:proj}
For $n\in\mathbb N$ let $\mathbf 1\in\R^n$ denote the all-ones vector. For $y\in\R^n$ define its mean
\[
\bar y:=\frac1n\langle y,\mathbf 1\rangle=\frac1n\sum_{i=1}^n y_i,
\]
and define the \emph{mean-zero projection} $P:\R^n\to\R^n$ by
\begin{equation}\label{eq:Pdef}
P(y):=y-\bar y\,\mathbf 1.
\end{equation}
equivalently, $P$ is the orthogonal projection onto the subspace $\{z\in\R^n:\langle z,\mathbf 1\rangle=0\}$.
\end{definition}

\begin{lemma}\label{lem:Pkernel}
For $y\in\R^n$ with mean $\bar y:=\frac1n\langle y,\mathbf 1\rangle$, the projection $P(y)=y-\bar y\,\mathbf 1$
vanishes if and only if $y$ is a constant vector, i.e. $y=\bar y\,\mathbf 1$.
In particular, if additionally $\langle y,\mathbf 1\rangle=0$, then $P(y)=0$ implies $y=0$.
\end{lemma}

\begin{proof}
By definition, $P(y)=0$ is equivalent to $y=\bar y\,\mathbf 1$, which is exactly the condition that all coordinates
of $y$ are equal to the constant $\bar y$.
If moreover $\langle y,\mathbf 1\rangle=0$, then $0=\langle y,\mathbf 1\rangle=\langle \bar y\,\mathbf 1,\mathbf 1\rangle=n\bar y$,
so $\bar y=0$ and hence $y=0$.
\end{proof}


\begin{definition}\label{def:defect}
For $x\in(\Rp)^n$ let $y=\log x$ and define the \emph{defect} of $x$ by
\[
\mathrm{Def}(x):=\norm{P(y)}_2.
\]
We say that $x$ has \emph{zero defect} if $\mathrm{Def}(x)=0$.
\end{definition}

\begin{remark}
Under the conservation constraint $\sigma(x)=\sum_i \log x_i=0$,
Lemma~\ref{lem:Pkernel} implies that $\mathrm{Def}(x)=0$ forces $\log x=0$ and hence $x\equiv 1$.
\end{remark}

\subsubsection{Coercivity: energy gap controls defect}\label{sec:coercivity} The coercivity step converts cost information into a quantitative defect bound. It rests on the
elementary inequality $\cosh(t)-1\geq t^2/2$. \begin{lemma}\label{lem:coshcoercive}
For all $t\in\R$,
\[
\cosh(t)-1\ \geq\ \frac{t^2}{2},
\]
with equality if and only if $t=0$.
\end{lemma}

\begin{proof}
Define $f(t):=\cosh(t)-1-\tfrac12 t^2$. Then $f(0)=0$, $f'(t)=\sinh(t)-t$ and
$f''(t)=\cosh(t)-1\geq 0$. Thus $f'$ is increasing with $f'(0)=0$, hence $f'(t)\geq 0$ for $t\geq 0$
and $f'(t)\leq 0$ for $t\le 0$, so $t=0$ is a global minimum of $f$. Therefore $f(t)\geq 0$ for all $t$,
with equality only at $t=0$.
\end{proof} \begin{theorem}\label{thm:coercive}
Let $x\in(\Rp)^n$ and $y=\log x$. Then
\begin{equation}\label{eq:coercive}
\sum_{i=1}^n J\!\left(e^{P(y)_i}\right)\ \geq\ \frac12\,\norm{P(y)}_2^2 \;=\;\frac12\,\mathrm{Def}(x)^2.
\end{equation}
In particular, if $\sum_{i=1}^n J\!\left(e^{P(y)_i}\right)=0$ then $\mathrm{Def}(x)=0$.
\end{theorem}

\begin{proof}
By Lemma~\ref{lem:logform}, $J(e^{P(y)_i})=\cosh(P(y)_i)-1$.
Apply Lemma~\ref{lem:coshcoercive} termwise and sum over $i$:
\[
\sum_i J(e^{P(y)_i})=\sum_i(\cosh(P(y)_i)-1)\ \geq\ \sum_i \frac{P(y)_i^2}{2}
=\frac12 \norm{P(y)}_2^2.
\]
If the left-hand side equals $0$, then each term equals $0$, hence each $P(y)_i=0$ and $\mathrm{Def}(x)=0$.
\end{proof}
% - END: added sections (v2) -

\subsection{\texorpdfstring{Finite-data aggregation: $W$-block window measurements}{Finite-data aggregation: W-block window measurements}}\label{sec:aggregation}

\begin{remark}\label{rem:Wparam}
All arguments in Sections~\ref{sec:coercivity}-\ref{subsec:domination} apply verbatim for any fixed
integer $W\geq 1$ once the corresponding measurement operator is specified.
In the examples we keep the concrete choice $W=8$.
\end{remark}

\subsection{\texorpdfstring{Window operator at resolution $W$}{Window operator at resolution W}}
Fix an integer $W\ge 1$ (e.g. $W=8$). For a real sequence $y=(y_0,y_1,\dots)$ define its \emph{$W$-block window sums}
\begin{equation}\label{eq:window}
\mathcal{W}(y)_k \;:=\; \sum_{j=0}^{W-1} y_{(Wk+j)},\qquad k=0,1,2,\dots.
\end{equation}
For brevity we sometimes write $W_k:=\mathcal W(y)_k$ for the window sums; this notation should not be confused with the fixed window \emph{length} $W$.
equivalently, $\mathcal{W}(y)$ is obtained by partitioning the time axis into consecutive blocks
of length $W$ and summing within each block (e.g. $W=8$).
\begin{remark}
The choice of fixed $W$ is treated as part of the finite-resolution model (C3); $W=8$ is used only as a concrete running example. The analysis requires only that the block length is a fixed positive integer.
\end{remark}

\subsection{The rational (finite-state) class and the degree parameter}
Let $y=(y_n)_{n\geq 0}$ be a real sequence with generating function
\[
\theta(z):=\sum_{n\geq 0} y_n z^n.
\]
We say that $y$ is \emph{rational of degree at most $d$} if $\theta(z)=p(z)/q(z)$ where
$\deg p\le d$, $\deg q\le d$, and $q(0)\neq 0$; we then normalise scale so that $q(0)=1$.
Write
\[
q(z)=1+q_1 z+q_2 z^2+\cdots+q_d z^d,
\]
allowing $q_d=0$ when $\deg q<d$ (padding to length $d$). In this setting, $d$ counts the number of
free recurrence coefficients (equivalently, the McMillan degree / state dimension of a minimal
linear realisation).

\begin{lemma}\label{lem:recurrence}
If $\theta(z)=p(z)/q(z)$ with $q(0)\neq 0$ (so a Taylor expansion at $0$ exists) and after normalisation $q(0)=1$ and $q(z)=1+q_1 z+\cdots+q_d z^d$, then the coefficients
$y_n$ satisfy the linear recurrence
\begin{equation}\label{eq:recurrence}
y_n + q_1 y_{n-1}+q_2 y_{n-2}+\cdots+q_d y_{n-d}=0,\qquad n\ge d+1.
\end{equation}
Conversely, any sequence satisfying \eqref{eq:recurrence} for some $(q_1,\dots,q_d)$ is determined by the
parameters $(q_1,\dots,q_d)$ and the initial values $(y_0,\dots,y_d)$, hence by $2d+1$ real numbers.
\end{lemma}

\begin{proof}

Since $\theta(z)=p(z)/q(z)$ and $q(0)=1$, we have the identity
\[
q(z)\,\theta(z)=p(z).
\]
Write $q(z)=1+q_1 z+\cdots+q_d z^d$ and $\theta(z)=\sum_{n\ge 0} y_n z^n$.
Comparing coefficients of $z^n$ for $n\ge d+1$ gives
\[
y_n+q_1y_{n-1}+\cdots+q_dy_{n-d}=0,
\]
because $\deg p\le d$ implies the $z^n$ coefficient of $p$ is $0$ for $n\ge d+1$.

Conversely, if \eqref{eq:recurrence} holds for fixed $(q_1,\dots,q_d)$, then each $y_n$ with $n\ge d+1$ is determined recursively from $(y_0,\dots,y_d)$. Hence the full sequence is determined by $(q_1,\dots,q_d,y_0,\dots,y_d)$, i.e.\ by $2d+1$ real parameters.

\end{proof}

\subsection{\texorpdfstring{Unique determination from $K\ge 2d+1$ window sums}{Unique determination from K>=2d+1 window sums}}\label{subsec:window-uniq}
We now fill in the self-contained ``unique determination from window data'' argument in a form
that does \emph{not} rely on external machine-verification tooling. Throughout this subsection we fix integers
\[
d\ge 1,\qquad W\ge 1,\qquad K\ge 2d+1.
\]
Let $\theta(z)=\sum_{n\ge 0}y_n z^n$ be a rational signal of degree $\le d$ with representation
$\theta(z)=p(z)/q(z)$ where $\deg p\le d$, $\deg q\le d$, and we normalise $q(0)=1$. Write
\[
q(z)=1+q_1 z+\cdots+q_d z^d.
\]
By Lemma~\ref{lem:recurrence}, once the parameter vector
\[
\pi:=(y_0,\dots,y_d,q_1,\dots,q_d)\in\R^{2d+1}
\]
is fixed, the recurrence \eqref{eq:recurrence} determines $y_n$ for all $n\ge d+1$.

Recall the $W$-block window sums $\mathcal W(y)_k$ from \eqref{eq:window}, and define the truncated window map (using the first $2d+1$ windows)
\[
F_{d,W}:\R^{2d+1}\to\R^{2d+1},\qquad
F_{d,W}(\pi):=\bigl(\mathcal W(y)_0,\dots,\mathcal W(y)_{2d}\bigr),
\]
where $y$ is generated from $\pi$ by \eqref{eq:recurrence}.
Here $D F_{d,W}(\pi)$ denotes the Jacobian matrix of $F_{d,W}$ at $\pi$.


\begin{lemma}\label{lem:window-generic-local}
The map $F_{d,W}$ is polynomial in $\pi$.
Indeed, each $y_n$ is obtained from $(y_0,\dots,y_d,q_1,\dots,q_d)$ by iterating the recurrence \eqref{eq:recurrence}, so $y_n$ is a polynomial function of $\pi$ for every fixed $n$; each window sum is a finite linear combination of such $y_n$, hence $F_{d,W}$ is polynomial. If there exists \emph{one} parameter point $\pi_0$
with $\det DF_{d,W}(\pi_0)\neq 0$, then
\[
\Omega_{d,W}:=\{\pi\in\R^{2d+1}:\det DF_{d,W}(\pi)\neq 0\}
\]
is a nonempty Zariski-open (hence open dense) subset of $\R^{2d+1}$, and on $\Omega_{d,W}$
the first $2d+1$ window sums determine $\pi$ \emph{locally uniquely} (i.e.\ $F_{d,W}$ is locally
one-to-one).
\end{lemma}

\begin{proof}
Each coefficient $y_n$ is obtained by iterating the linear recurrence \eqref{eq:recurrence} whose
update rule is polynomial in $(q_1,\dots,q_d)$ and linear in $(y_0,\dots,y_d)$. Hence each
$\mathcal W(y)_k$ is polynomial in $\pi$, so $F_{d,W}$ is polynomial and $\det DF_{d,W}$ is a
polynomial function. If $\det DF_{d,W}(\pi_0)\neq 0$ for some $\pi_0$, then the polynomial
$\det DF_{d,W}$ is not identically zero, so its nonvanishing locus $\Omega_{d,W}$ is Zariski-open
and nonempty. For $\pi\in\Omega_{d,W}$, the inverse function theorem gives local invertibility,
hence local uniqueness of $\pi$ from the first $2d+1$ windows.
\end{proof}



\begin{theorem}\label{thm:generic-full-rank}
Fix integers $d\ge 1$ and $W\ge 1$. Assume there exists a parameter point $\pi_0$ such that $\det DF_{d,W}(\pi_0)\neq 0$. Then the nonvanishing locus
\[
\Omega_{d,W}:=\{\pi:\det DF_{d,W}(\pi)\neq 0\}
\]
is nonempty and Zariski-open (hence open dense) in $\R^{2d+1}$, and for every $\pi\in\Omega_{d,W}$ the window map $F_{d,W}$ is locally invertible at $\pi$. Moreover, to certify the hypothesis it suffices to exhibit an integer point $\pi_0\in\Z^{2d+1}$ and a prime $p$ for which $\det DF_{d,W}(\pi_0)\not\equiv 0\pmod p$.
\end{theorem}

\begin{proof}
By Lemma~\ref{lem:window-generic-local}, $\det DF_{d,W}$ is a polynomial in the coordinates of $\pi$. If $\det DF_{d,W}(\pi_0)\neq 0$, then this polynomial is not identically zero, so $\Omega_{d,W}$ is Zariski-open and nonempty; the inverse function theorem yields local invertibility at each $\pi\in\Omega_{d,W}$. For the finite-field certificate: if $\pi_0\in\Z^{2d+1}$ and $\det DF_{d,W}(\pi_0)\not\equiv 0\pmod p$, then the integer $\det DF_{d,W}(\pi_0)$ is not divisible by $p$, hence is nonzero over $\Z$, and therefore nonzero as a real number.
\end{proof}

\begin{remark}
Theorem~\ref{thm:generic-full-rank} reduces generic full-rank (and hence local identifiability) for a given $(d,W)$ to finding \emph{one} explicit witness point. In practice, such witnesses can often be found by a brief randomized integer search, and then certified by a modular determinant computation as above.
\end{remark}

\begin{lemma}\label{lem:witness-38}
For $d=3$ and $W=8$, the set $\Omega_{3,8}$ from Lemma~\ref{lem:window-generic-local} is nonempty. In particular, this verifies the hypothesis of Theorem~\ref{thm:generic-full-rank} for $(d,W)=(3,8)$.
In particular, for the parameter point
\[
\pi_0=(y_0,y_1,y_2,y_3,q_1,q_2,q_3)=(1,1,5,1,2,2,-2)
\]
one has $\det DF_{3,8}(\pi_0)\neq 0$.
Consequently $\Omega_{3,8}$ is Zariski-open (hence open dense) in $\R^{7}$, and the first $2d+1=7$
window sums determine $\pi$ locally uniquely throughout $\Omega_{3,8}$.
\end{lemma}

\begin{proof}
We certify $\det DF_{3,8}(\pi_0)\neq 0$ by a finite check in a prime field.
Let $p:=10^9+7$ and work in $\mathbb{F}_p$ throughout.
Using the recurrence \eqref{eq:recurrence} at $\pi=\pi_0$, compute $y_n$ up to $n=55$ and hence the first $K=7$ window sums
\[
\begin{aligned}
(W_0,\dots,W_6)=(&-16,-5160,-975168,-169890432,\\
&-27959752704,-4399334572032,-665805326548992).
\end{aligned}
\]

Next propagate parameter derivatives through the same recurrence:
for each parameter $\alpha\in\{y_0,y_1,y_2,y_3,q_1,q_2,q_3\}$ set
$u^{(\alpha)}_n:=\partial y_n/\partial \alpha$ and differentiate
$y_n+q_1y_{n-1}+q_2y_{n-2}+q_3y_{n-3}=0$ $(n\ge 4)$ to obtain
\[
u^{(\alpha)}_n+q_1u^{(\alpha)}_{n-1}+q_2u^{(\alpha)}_{n-2}+q_3u^{(\alpha)}_{n-3}
=-(\partial_\alpha q_1)y_{n-1}-(\partial_\alpha q_2)y_{n-2}-(\partial_\alpha q_3)y_{n-3}.
\]
These recurrences determine all $u^{(\alpha)}_n$ needed for the window derivatives
$\partial_\alpha W_k=\sum_{j=0}^{7}u^{(\alpha)}_{8k+j}$.
Assemble the $7\times 7$ Jacobian $DF_{3,8}(\pi_0)$ with entries $\partial_\alpha W_k$.
(The full integer Jacobian matrix is recorded in Appendix~\ref{app:jacobian-38}; here we retain only the determinant residue modulo $p$.)

Compute the determinant in $\mathbb{F}_p$ by Gaussian elimination. The result is
\begin{equation*}
\det\bigl(DF_{3,8}(\pi_0)\bigr)\equiv 972226939\pmod{p}.
\end{equation*}
, which is nonzero in $\mathbb{F}_p$. Hence $\det DF_{3,8}(\pi_0)\neq 0$ over the integers.
\end{proof}


\begin{remark}\label{rem:witness-points}
For any fixed pair $(d,W)$, the determinant $\det DF_{d,W}(\pi)$ is a polynomial function of the
parameter vector $\pi$ (with coefficients depending only on $(d,W)$). Hence the full-rank locus
\[
\Omega_{d,W}:=\{\pi:\det DF_{d,W}(\pi)\neq 0\}
\]
is either empty or Zariski open (and therefore Euclidean open and dense) in the ambient parameter space.
To certify nonemptiness in the present paper (working entirely within the present manuscript),
it suffices to exhibit a single explicit witness $\pi_0$ with $\det DF_{d,W}(\pi_0)\neq 0$.
Lemma~\ref{lem:witness-38} provides such a witness for the operating example $(d,W)=(3,8)$.
Additional witnesses for other $(d,W)$ pairs can be recorded in the same way.
\end{remark}

\begin{table}[t]
	\centering
	\begin{tabular}{lll}
		\hline
		$(d,W)$ & witness $\pi_0$ & certificate of full rank \\
		\hline \\
		$(3,8)$ & $(1,1,5,1,2,2,-2)$ & Lemma~\ref{lem:witness-38} \\
		\hline
	\end{tabular}
	\begingroup
	\captionsetup{font=footnotesize}
	\caption{Example full-rank witness points. Further rows can be added by exhibiting any $\pi_0$ with $\det DF_{d,W}(\pi_0)\neq 0$.}
	\endgroup
\end{table}





\begin{theorem}\label{thm:window-ident-general}
Assume $\Omega_{d,W}$ is nonempty (equivalently, $\det DF_{d,W}$ is not the zero polynomial).
Then for every $\pi\in\Omega_{d,W}$ the first $K$ window sums $\bigl(\mathcal W(y)_0,\dots,\mathcal W(y)_{K-1}\bigr)$
determine $\pi$ locally uniquely within the degree-$\le d$ rational class. In particular, the
degree-$\le d$ signal $\theta$ and the full coefficient sequence $(y_n)_{n\ge 0}$ are locally
uniquely determined on $\Omega_{d,W}$ by $K\ge 2d+1$ windows.
\end{theorem}

\begin{proof}
Fix $\pi\in\Omega_{d,W}$. By Lemma~\ref{lem:window-generic-local} the map
\[
F_{d,W}:\R^{2d+1}\to\R^{2d+1},\qquad
F_{d,W}(\pi)=\bigl(\mathcal W(y)_0,\dots,\mathcal W(y)_{2d}\bigr)
\]
has nonsingular Jacobian at $\pi$, hence is locally injective near $\pi$ (Inverse Function Theorem).
Now let $K\ge 2d+1$ and define
\[
G_{K}(\pi):=\bigl(\mathcal W(y)_0,\dots,\mathcal W(y)_{K-1}\bigr)\in\R^{K}.
\]
Let $\mathrm{pr}:\R^{K}\to\R^{2d+1}$ be the coordinate projection onto the first $2d+1$ components.
Then $\mathrm{pr}\circ G_{K}=F_{d,W}$. Therefore, if $G_{K}(\pi')=G_{K}(\pi)$ for some $\pi'$,
then $F_{d,W}(\pi')=F_{d,W}(\pi)$. Taking $\pi'$ sufficiently close to $\pi$, local injectivity of
$F_{d,W}$ forces $\pi'=\pi$. Hence the first $K$ windows determine $\pi$ locally uniquely on
$\Omega_{d,W}$ for every $K\ge 2d+1$.
\end{proof}

\begin{remark}\label{rem:window-global}
Theorem~\ref{thm:window-ident-general} is the strongest statement that follows from a Jacobian
(non-singularity) argument alone: it guarantees \emph{local} identifiability on an open dense set.
A \emph{global} uniqueness theorem (one-to-one on a parameter locus) requires additional structure or
additional observables beyond window sums to rule out distinct far-apart preimages with the same window data.
\end{remark}

\begin{theorem}\label{thm:window-ident-global}
Let $S_k:=\mathcal W(y)_k$ be the window-sum sequence for a degree-$\le d$ rational signal.
Assume that $S_k$ admits an exponential-sum representation
\[
S_k=\sum_{i=1}^d A_i\,\mu_i^k
\]
\noindent\textbf{(H-prony)}\; Assume the exponents $\mu_1,\dots,\mu_d$ are distinct and nonzero, the amplitudes $A_i$ are nonzero, and the $d\times d$ Hankel matrix $H:=(S_{i+j})_{i,j=0}^{d-1}$ is invertible. (This is a standard nondegeneracy / local invertibility condition for Prony-type reconstruction.)
Then the data $(S_0,\dots,S_{2d-1})$ determine uniquely:
\begin{enumerate}
\item the monic degree-$d$ characteristic polynomial $\chi(t)=\prod_{i=1}^d(t-\mu_i)$ of the order-$d$ recurrence satisfied by $(S_k)$, hence the multiset $\{\mu_i\}_{i=1}^d$;
\item the amplitudes $(A_i)_{i=1}^d$, hence the entire sequence $(S_k)_{k\ge 0}$.
\end{enumerate}
Equivalently, the collection of parameter pairs $(\mu_i,A_i)_{i=1}^d$ is determined up to a common permutation of indices. If one fixes any deterministic ordering convention on the distinct exponents (for example, ordering by $|\mu_i|$ when these moduli are pairwise distinct), then the ordered tuple $(\mu_1,\dots,\mu_d,A_1,\dots,A_d)$ is uniquely determined.

Consequently, for $K\ge 2d$ the first $K$ window sums determine the induced window-sum process uniquely on this nondegenerate locus.
\end{theorem}

\begin{proof}
The coefficients of $P(t)=t^d+a_1 t^{d-1}+\cdots+a_d$ solve the Hankel linear system
\[
\sum_{m=1}^d a_m\,S_{k+d-m}=-S_{k+d}\qquad (k=0,1,\dots,d-1),
\]
whose matrix is precisely $H=(S_{i+j})_{i,j=0}^{d-1}$. By hypothesis $H$ is invertible, hence
$(a_1,\dots,a_d)$ (and thus $P$) are uniquely determined from $(S_0,\dots,S_{2d-1})$.
Since the $\mu_i$ are distinct, the roots of $\chi$ are exactly $\{\mu_i\}_{i=1}^d$.
With $\mu_i$ known, the amplitudes $(A_i)$ are obtained uniquely from the Vandermonde system
\[
(S_0,\dots,S_{d-1})^\top = V(\mu)\,(A_1,\dots,A_d)^\top,
\quad V(\mu)_{k,i}=\mu_i^{\,k},
\]
whose determinant is $\prod_{i<j}(\mu_j-\mu_i)\neq 0$. This determines $(S_k)$ for all $k$.
\end{proof}


\begin{definition}[Local optimality on an identifiability locus]\label{def:local-optimal}
Fix a neighborhood $\mathcal{U}\subset\mathcal{M}_{d,W}$.
A sound procedure $\Psi$ is \emph{locally optimal on $\mathcal{U}$} if for every sound procedure $\Xi$ one has
\[
\operatorname{Res}(\Xi)\cap\mathcal{U}\subseteq \operatorname{Res}(\Psi)\cap\mathcal{U}
\]
and the outputs agree on every instance resolved by $\Xi$ in $\mathcal{U}$, i.e.\ for all $y\in \operatorname{Res}(\Xi)\cap\mathcal{U}$,
\[
\Psi(\mathcal{W}_{W,K}(y))=\Xi(\mathcal{W}_{W,K}(y)).
\]
\end{definition}

\begin{corollary}[Uniqueness among locally optimal rules]\label{cor:unique-local-optimal}
Under the hypotheses of Theorem~\ref{thm:domination}, the canonical procedure $\Phi^\ast$ is locally optimal on $\mathcal{U}_{d,W}$.
Moreover, if $\Psi$ is any locally optimal sound procedure on $\mathcal{U}_{d,W}$, then
\[
\operatorname{Res}(\Psi)\cap\mathcal{U}_{d,W}=\operatorname{Res}(\Phi^\ast)\cap\mathcal{U}_{d,W},
\]
and for every $y\in \operatorname{Res}(\Phi^\ast)\cap\mathcal{U}_{d,W}$ the outputs agree:
\[
\Psi(\mathcal{W}_{W,K}(y))=\Phi^\ast(\mathcal{W}_{W,K}(y)).
\]
\end{corollary}

\begin{proof}
By Theorem~\ref{thm:domination}, $\Phi^\ast$ dominates every sound procedure on $\mathcal{U}_{d,W}$, hence it is locally optimal by Definition~\ref{def:local-optimal}.
If $\Psi$ is locally optimal, then applying Definition~\ref{def:local-optimal} with $\Xi=\Phi^\ast$ gives
$\operatorname{Res}(\Phi^\ast)\cap\mathcal{U}_{d,W}\subseteq \operatorname{Res}(\Psi)\cap\mathcal{U}_{d,W}$ and agreement on that set.
Together with Theorem~\ref{thm:domination} applied to $\Psi$ (which yields the reverse inclusion and agreement), this gives the claimed equality of resolve sets and agreement of outputs.
\end{proof}



\begin{remark}[Parameter uniqueness up to permutation / with ordering]
Theorem~\ref{thm:window-ident-global} determines the unordered node set
$\{\mu_1,\dots,\mu_d\}$ uniquely, hence parameters are unique up to simultaneous permutation of pairs
$(\mu_i,A_i)$.
If one imposes an ordering convention (for example $\mu_1<\cdots<\mu_d$ when nodes are real and positive),
the parameter tuple itself is uniquely determined.
\end{remark}


\begin{remark}\label{rem:generic-uniqueness}
In concrete Prony/Hankel implementations, failure of the hypotheses in Theorem~\ref{thm:window-ident-global}
(distinct exponents, nonzero amplitudes, and invertible Hankel/Vandermonde matrices) occurs on a proper algebraic subset
of parameter space (e.g. discriminant and determinant vanishing conditions). Hence the stated nondegeneracy conditions
are \emph{generic} (open dense) within standard finite-dimensional parametrizations of degree-$\le d$ rational signals.
\end{remark}


\begin{theorem}[General rank certification]\label{thm:general-rank}
For every pair of integers $d\ge 1$ and $W\ge 1$, the identifiability locus
$\Omega_{d,W}$ is nonempty.
In particular, the generic local uniqueness conclusion of Theorem~\ref{thm:window-ident-general}
and the global Prony recovery of Theorem~\ref{thm:window-ident-global} apply for every $(d,W)$ under the corresponding nondegeneracy hypotheses.
\end{theorem}

\begin{proof}
We exhibit a witness family.
Choose distinct $\alpha_1,\dots,\alpha_d\in(0,1)$ (e.g.\ $\alpha_i:=1/(i+1)$), and define
$y_n:=\sum_{i=1}^d \alpha_i^n$.
Its $W$-block window sums satisfy
\[
S_k=\sum_{i=1}^d B_i\,\mu_i^k,\qquad
\mu_i:=\alpha_i^W,\quad
B_i:=\frac{1-\alpha_i^W}{1-\alpha_i}.
\]
The nodes $\mu_i$ are distinct and $B_i>0$.
Hence the Hankel matrix $H=(S_{i+j})_{0\le i,j\le d-1}$ admits
\[
H=V^\top\!\operatorname{diag}(B)\,V,
\]
with $V$ the Vandermonde matrix on $(\mu_i)$, so
\[
\det H=(\det V)^2\prod_{i=1}^d B_i
=\left(\prod_{i<j}(\mu_j-\mu_i)\right)^2\prod_{i=1}^d B_i\neq 0.
\]
Therefore the Prony/Hankel nondegeneracy conditions hold at this witness.
By Lemma~\ref{lem:window-generic-local}, this gives a nonempty full-rank identifiability locus.
\end{proof}


\begin{corollary}\label{cor:window-ident-38}
In the specialized regime used in this paper we take $d=3$ and $W=8$, hence $K_{\min}=2d+1=7$.
We emphasize the two different thresholds: $K\ge 2d+1$ is the dimensional requirement for \emph{local} recovery of the full parameter $\pi\in\R^{2d+1}$ from the truncated window map $F_{d,W}$, while $K\ge 2d$ is sufficient for \emph{global} recovery of the induced window-sum process via Prony/Hankel methods under the nondegeneracy hypotheses of Theorem~\ref{thm:window-ident-global}.
For this choice, Theorem~\ref{thm:general-rank} yields $\Omega_{3,8}\neq\emptyset$, so the generic local uniqueness conclusion of
Theorem~\ref{thm:window-ident-general} holds. Moreover, under the nondegeneracy hypotheses of Theorem~\ref{thm:window-ident-global}, the first $2d$ window sums (hence any $K\ge 2d$) determine the window-sum sequence (equivalently, the exponential parameters $(\mu_i,A_i)$) \emph{globally} uniquely.
\end{corollary}

\begin{proof}
The specialization $d=3$, $W=8$ fixes $K_{\min}=2d+1=7$ and the local threshold statement is the
instance of Theorem~\ref{thm:window-ident-general} on the nonempty full-rank locus $\Omega_{3,8}$
certified in Lemma~\ref{lem:witness-38}.
The global $K\ge 2d$ threshold for recovery of the induced window-sum process (equivalently, the exponential parameters)
follows from Theorem~\ref{thm:window-ident-global} under its stated Prony/Hankel nondegeneracy hypotheses.
\end{proof}


\section{Certification and optimality from finite-window data}
\label{sec:cert-opt}
This section develops the finite-data certification pipeline built on the window operator:
we pass from raw window measurements to a canonical defect score, then to a decision rule.
We isolate the status of the window length $W$, give $\varepsilon$-robust certification bounds,
and define soundness and domination among finite-data procedures, culminating in
$\varepsilon$-optimal certification statements.



\begin{theorem}[Master Theorem (architectural summary)]\label{thm:master}
Fix integers $d\ge 0$, $W\ge 1$, and $K\ge 2d+1$.  Consider the window map $F_{d,W}$ and the finite-data pipeline
\[
\Phi^\ast:\ w\ \mapsto\ A(w)\ \mapsto\ x\ \mapsto\ u=P(\log x)\ \mapsto\ B(u)\in\{\zero,\nonzero,\inconclusive\},
\]
as defined in Section~\ref{sec:cert-opt}.  Assume the standing hypotheses needed for the cited results below
(in particular, the cost axioms and the identifiability assumptions for the window model).  Then:
\begin{enumerate}[label=\textup{(\roman*)},leftmargin=2.2em]
\item \textbf{Soundness.}  The decision rule $\Phi^\ast$ is sound in the sense of Definition~\ref{def:sound}.
\item \textbf{$\varepsilon$-robust certification.}  Under perturbations of size $\varepsilon$, the certified conclusion degrades in the sharp two-sided way, yielding membership in $\Mean_{2\varepsilon}$ as in Theorems~\ref{thm:eps-tolerant} and~\ref{thm:eps-cert}.
\item \textbf{Local domination (optimality).}  On the identifiability neighborhood $\mathcal U_{d,W}$, $\Phi^\ast$ dominates every sound procedure (Theorem~\ref{thm:domination}).
\item \textbf{Uniqueness among locally optimal rules.}  Any sound procedure that is locally optimal on $\mathcal U_{d,W}$ agrees with $\Phi^\ast$ on all resolved data in $\mathcal U_{d,W}$ (Corollary~\ref{cor:unique-local-optimal}).
\item \textbf{Forced decision-layer dependence.}  Any sound defect-based decision rule factors through the scale-invariant invariant $u=P(\log x)$, hence through the $(P,B)$-core in the sense of Remark~\ref{rem:forced-PB}.
\item \textbf{Non-vacuity via generic full rank.}  For any $(d,W)$, if there exists a witness $\pi_0$ with $\det DF_{d,W}(\pi_0)\neq 0$, then the full-rank locus is a nonempty Zariski-open (hence dense) subset (Theorem~\ref{thm:generic-full-rank}); for $(d,W)=(3,8)$ such a witness is certified in Lemma~\ref{lem:witness-38}.
\end{enumerate}
If, in addition, the window map is globally injective on the parameter locus $\mathcal M_{d,W}$, then the domination statement upgrades from $\mathcal U_{d,W}$ to $\mathcal M_{d,W}$ (Corollary~\ref{cor:domination-global}).
\end{theorem}

\begin{proof}
Items (i)–(ii) follow from Definition~\ref{def:sound} together with Theorems~\ref{thm:eps-tolerant} and~\ref{thm:eps-cert}.  Item (iii) is Theorem~\ref{thm:domination}.  Item (iv) is Corollary~\ref{cor:unique-local-optimal}.  Item (v) is the factorization/forced-dependence principle recorded in Remark~\ref{rem:forced-PB}.  Item (vi) is Theorem~\ref{thm:generic-full-rank}, with nonemptiness witnessed for $(3,8)$ by Lemma~\ref{lem:witness-38}.  The final global upgrade is Corollary~\ref{cor:domination-global}.
\end{proof}


\subsection{Status of the window length \texorpdfstring{$W$}{W}}\label{subsec:status-W}
Throughout, $W\in\N$ denotes the fixed block length used to form window sums \eqref{eq:window}.
The theory is stated for general $W$; the value $W=8$ is a modeling choice, used below only as an
illustrative operating point.

A simple (self-contained) motivation for the number $8$ comes from a three-bit latent state
$b\in\{0,1\}^3$. A Gray-code cycle on the 3-cube visits all $2^3=8$ states exactly once before
returning to the start, for example
\[
000\to 001\to 011\to 010\to 110\to 111\to 101\to 100\to 000,
\]
where successive states differ in exactly one bit. If a measurement/aggregation step is designed to
cover one full cycle of such a three-bit state evolution, then the natural ``cycle window'' length is
$W=8$. This is only a heuristic design rationale: none of the main theorems depends on it, and the
identifiability threshold remains $K\ge 2d+1$ for any fixed $W\ge 1$.

\subsection{\texorpdfstring{$\varepsilon$}{eps}-tolerant certification from noisy window sums}\label{subsec:eps-stability}
In practice, window sums are observed with finite precision. Instead of exact constraints
$\mathcal{W}(y)_k=0$, we assume a uniform tolerance model
\begin{equation}\label{eq:window-noise}
\bigl|W_k\bigr|\le \varepsilon,\qquad k=0,1,\dots,K-1.
\end{equation}
The next result shows that small window-sum errors yield a small certification cost for the
reconstructed signal. The bound is quadratic in $\varepsilon$, with constants controlled by the
conditioning of the reconstruction map.

\begin{lemma}\label{lem:cosh-upper}
For every $t\in\R$,
\[
\cosh(t)-1 \ \le\ \frac12\,e^{|t|}\,t^2.
\]
In particular, if $|t|\le \tau$ then $\cosh(t)-1\le \frac12 e^{\tau}t^2$.
\end{lemma}

\begin{proof}
Using the power series $\cosh(t)-1=\sum_{m\ge 1}\frac{t^{2m}}{(2m)!}$ and $(2m)!\ge 2\,(2m-2)!$,
\[
\cosh(t)-1 \le \frac{t^2}{2}\sum_{m\ge 1}\frac{|t|^{2m-2}}{(2m-2)!}
= \frac{t^2}{2}\cosh(|t|)\le \frac{t^2}{2}e^{|t|}.
\]
\end{proof}


\begin{proposition}\label{prop:rec-exists}
Fix $W\in\N$, $d\ge 0$, and $n\in\N$, and set $K:=\lceil n/W\rceil$.
Assume $K\ge 2d+1$ and work on a nondegenerate parameter locus where the truncated window map
\[
\pi\longmapsto \mathcal{W}_{W,K}(y(\pi))
\]
is locally invertible at the neutral realization (equivalently, its Jacobian determinant is nonzero).
Then there exist $\varepsilon_0>0$ and a $C^1$ reconstruction map
\[
\mathsf{Rec}_{d,W}:\{w\in\R^K:\ \|w\|_\infty\le \varepsilon_0\}\to\R^n
\]
such that $\mathsf{Rec}_{d,W}(w)$ returns the first $n$ coefficients of the unique model signal
compatible with the window data $w$ in that neighborhood. Moreover, writing
\[
L\:=\ \sup_{\|w\|_\infty\le \varepsilon_0}\ \bigl\|D\mathsf{Rec}_{d,W}(w)\bigr\|_{2\to 2},
\]
one has the Lipschitz bound
\[
\|\mathsf{Rec}_{d,W}(w)-\mathsf{Rec}_{d,W}(0)\|_2\le L\,\|w\|_2\qquad(\|w\|_\infty\le \varepsilon_0).
\]
\end{proposition}

\begin{proof}
Write $\|\cdot\|_\infty$ for the max norm on $\R^K$, i.e.\ $\|w\|_\infty:=\max_{0\le k\le K-1}|w_k|$, and $\|\cdot\|_2$ for the Euclidean norm.
Local invertibility of $\pi\mapsto \mathcal{W}_{W,K}(y(\pi))$ at the neutral realization implies, by the inverse function theorem, that there exist
neighborhoods $U$ of the neutral parameter and $V$ of $w=0$ such that the window map restricts to a $C^1$ diffeomorphism $U\to V$.
Define $\mathsf{Rec}_{d,W}$ on $V$ by taking the inverse map $V\to U$ and then returning the first $n$ coefficients of the corresponding model signal;
shrinking $V$ if necessary yields the stated cube $\{w:\|w\|_\infty\le \varepsilon_0\}\subseteq V$.

For the Lipschitz bound, $\mathsf{Rec}_{d,W}$ is $C^1$ on the closed cube, hence its derivative operator norm
$L:=\sup_{\|w\|_\infty\le \varepsilon_0}\|D\mathsf{Rec}_{d,W}(w)\|_{2\to 2}$ is finite.
By the mean value inequality on line segments,
\[
\|\mathsf{Rec}_{d,W}(w)-\mathsf{Rec}_{d,W}(0)\|_2
\le \sup_{t\in[0,1]}\|D\mathsf{Rec}_{d,W}(tw)\|_{2\to 2}\,\|w\|_2
\le L\,\|w\|_2,
\]
as claimed.
\end{proof}


\begin{theorem}\label{thm:eps-tolerant}
Fix $W\in\N$, $d\ge 0$, and $n\in\N$, and set $K:=\lceil n/W\rceil$.
Assume $K\ge 2d+1$ and that Proposition~\ref{prop:rec-exists} holds with constants $(\varepsilon_0,L)$.
Given window data $w\in\R^K$ with $|w_k|\le \varepsilon\le \varepsilon_0$, let $\hat y:=\mathsf{Rec}_{d,W}(w)$.
Assume the noiseless reconstruction $\hat y^{(0)}:=\mathsf{Rec}_{d,W}(0)$ is neutral, i.e.\ $P(\hat y^{(0)})=0$.
Then the certification objective
\[
\Phi(\hat y)\:=\ \sum_{i=1}^n J\!\left(e^{P(\hat y)_i}\right)\ =\ \sum_{i=1}^n \bigl(\cosh(P(\hat y)_i)-1\bigr)
\]
satisfies the explicit quadratic bound
\begin{equation}\label{eq:eps-bound}
\Phi(\hat y)\ \le\ \frac12\,e^{\,L\sqrt{K}\,\varepsilon_0}\,L^2K\,\varepsilon^2.
\end{equation}
\end{theorem}

\begin{proof}
Let $\hat y^{(0)}:=\mathsf{Rec}_{d,W}(0)$. By Proposition~\ref{prop:rec-exists}, for $\|w\|_\infty\le\varepsilon_0$ we have
\[
\|\hat y-\hat y^{(0)}\|_2\ \le\ L\,\|w\|_2 \ \le\ L\sqrt{K}\,\varepsilon.
\]
Since $P$ is an orthogonal projection, $\|P(u)-P(v)\|_2\le \|u-v\|_2$, hence using $P(\hat y^{(0)})=0$,
\[
\|P(\hat y)\|_2 \ =\ \|P(\hat y)-P(\hat y^{(0)})\|_2 \ \le\ \|\hat y-\hat y^{(0)}\|_2
\ \le\ L\sqrt{K}\,\varepsilon.
\]
Also $\|P(\hat y)\|_\infty\le \|P(\hat y)\|_2\le L\sqrt{K}\,\varepsilon \le L\sqrt{K}\,\varepsilon_0$.
Applying Lemma~\ref{lem:cosh-upper} componentwise gives
\[
\Phi(\hat y)=\sum_{i=1}^n(\cosh(P(\hat y)_i)-1)
\le \frac12 e^{\|P(\hat y)\|_\infty}\,\|P(\hat y)\|_2^2
\le \frac12 e^{L\sqrt{K}\varepsilon_0}\,L^2K\,\varepsilon^2,
\]
which is \eqref{eq:eps-bound}.
\end{proof}

\begin{remark}\label{rem:eps-numeric}
In the special case $(d,W)=(3,8)$ one may take $K=7$ windows.
Once a bound $L\le L_0$ is established for the local reconstruction map on $\|w\|_\infty\le \varepsilon_0$,
Theorem~\ref{thm:eps-tolerant} yields the explicit estimate
\[
\Phi(\hat y)\ \le\ \tfrac12 e^{L_0\sqrt{7}\,\varepsilon_0}\,L_0^2\cdot 7\,\varepsilon^2.
\]
For example, if a numerical conditioning check verifies $L_0\le 10$ on $\varepsilon_0=10^{-2}$,
then $\Phi(\hat y)\le 350\,\varepsilon^2$ for all $\varepsilon\le 10^{-2}$.
\end{remark}





\begin{remark}\label{rem:conditioning}
The Lipschitz constant $L$ in Proposition~\ref{prop:rec-exists} can be expressed in terms of the operator norm of
the inverse (or pseudoinverse) of the linearised reconstruction matrix at $w=0$; equivalently, in
terms of a condition number $\kappa$ for the Vandermonde/Hankel system used by the reconstruction.
Poor conditioning increases $L$ (equivalently the relevant condition number), but the dependence on $\varepsilon$ remains quadratic.
\end{remark}

\subsection{The certification objective}

\begin{definition}\label{def:meaning}
Given a cost function $c:S\times O\to\R_{\ge 0}$, the \emph{meaning map} (argmin correspondence) is
\[\Mean(s):=\arg\min_{o\in O} c(s,o):=\{o\in O: c(s,o)=\inf_{o'\in O}c(s,o')\}.\]
\end{definition}

\begin{remark}\label{rem:attainment}
The meaning map $\Mean(s)$ is defined via an infimum over $O$ and may be empty in full generality.
In the intended finite-data setting one works with a \emph{finite} candidate list $O$, in which case
every function $O\to\R$ attains its minimum and $\Mean(s)\neq\emptyset$ for all $s$.
Whenever we speak of ``the'' optimizer or of certificates that ``attain their minimum exactly on $\Mean(s)$'',
we implicitly restrict to such settings (or more generally assume that $c(s,\cdot)$ attains its minimum on $O$).
\end{remark}


\begin{proposition}[Attainment on compact candidate classes]\label{prop:meaning-compact}
Fix $s\in S$.
Assume $O$ is compact and the map $o\mapsto c(s,o)$ is continuous.
Then $\Mean(s)$ is nonempty.
\end{proposition}

\begin{proof}
By continuity on a compact set, $c(s,\cdot)$ attains its minimum at some $o_\ast\in O$
(extreme value theorem).
Hence
\[
c(s,o_\ast)=\inf_{o\in O}c(s,o),
\]
so $o_\ast\in\Mean(s)$ by Definition~\ref{def:meaning}.
\end{proof}


In many applications $O$ is not finite but is a closed subset of a Euclidean space; in that case one can ensure attainment by a standard coercivity hypothesis.


\begin{proposition}[Attainment under coercivity]\label{prop:meaning-coercive}
Fix $s\in S$ and assume $O\subset\R^m$ is nonempty and closed for some $m\ge 1$.
Assume $o\mapsto c(s,o)$ is lower semicontinuous on $O$ and \emph{coercive} on $O$ in the sense that
\[
c(s,o)\to\infty \quad\text{whenever}\quad \norm{o}\to\infty\ \text{ with }o\in O.
\]
Then $\Mean(s)$ is nonempty.
\end{proposition}

\begin{proof}
Let $\alpha:=\inf_{o\in O}c(s,o)$ and choose a minimizing sequence $(o_k)\subset O$ with $c(s,o_k)\downarrow\alpha$.
By coercivity, the sequence $(o_k)$ is bounded; thus it lies in $O\cap \overline{B_R}$ for some $R$.
By compactness of $\overline{B_R}$, there is a convergent subsequence $o_{k_j}\to o_\ast$.
Since $O$ is closed, $o_\ast\in O$.
Lower semicontinuity gives $c(s,o_\ast)\le \liminf_{j\to\infty} c(s,o_{k_j})=\alpha$, hence $c(s,o_\ast)=\alpha$ and $o_\ast\in\Mean(s)$.
\end{proof}


\begin{definition}\label{def:sound}
Fix $s\in S$ and write $c_s(o):=c(s,o)$ and $\mathcal{C}_s(o):=\mathcal{C}(s,o)$.
Assume throughout that the minimizer sets are nonempty (e.g.\ $O$ finite, or by Proposition~\ref{prop:meaning-compact} or Proposition~\ref{prop:meaning-coercive}).
A certificate $\mathcal{C}$ is
\begin{enumerate}[label=\textup{(\roman*)},leftmargin=2.2em]
\item \emph{minimizer-sound} if $\arg\min_{o\in O}\mathcal{C}_s(o)=\arg\min_{o\in O}c_s(o)$ for every $s$;
\item \emph{ordinally sound} if for every $s$ and all $o_1,o_2\in O$,
\[
\mathcal{C}_s(o_1)\le \mathcal{C}_s(o_2)\quad\Longleftrightarrow\quad c_s(o_1)\le c_s(o_2).
\]
\end{enumerate}
Ordinal soundness is strictly stronger than minimizer-soundness; it asserts that $\mathcal{C}_s$ induces
exactly the same preorder on candidates as the true cost $c_s$.
\end{definition}

\subsection{Sound procedures and domination}\label{subsec:domination}

We formulate the finite-data decision task as a three-valued procedure based on window aggregates.
Fix integers $W\ge 1$ and $K\ge 1$. For a signal $y=(y_n)_{n\ge 0}$ define the $W$-block window-sum
vector
\[
\mathcal{W}_{W,K}(y):=\big(\mathcal{W}(y)_0,\dots,\mathcal{W}(y)_{K-1}\big)\in\R^{K},
\qquad
\mathcal{W}(y)_k:=\sum_{j=0}^{W-1}y_{(Wk+j)}.
\]
We write the corresponding set of realizable window vectors as
\[
\mathcal{W}_{W,K}(\mathcal{M}_{d,W}):=\{\mathcal{W}_{W,K}(y):y\in\mathcal{M}_{d,W}\}\subset\R^{K}.
\]
A \emph{finite-data procedure} is a map
\[
\Psi:\R^{K}\to\{\zero,\nonzero,\inconclusive\}.
\]

\begin{definition}\label{def:domination}
Fix a model class $\mathcal{M}_{d,W}$ of signals (e.g.\ rational/finite-state of degree $\le d$) together with the
nondegeneracy hypotheses from Theorem~\ref{thm:window-ident-global}.
Let $y_{\mathrm{neu}}\in\mathcal{M}_{d,W}$ denote the distinguished neutral signal (one may take $y_{\mathrm{neu}}\equiv 0$).
We compare procedures only on realizable window vectors $w\in\mathcal{W}_{W,K}(\mathcal{M}_{d,W})$; outside this set a procedure may return $\inconclusive$ arbitrarily.
A procedure $\Psi$ is \emph{sound on $\mathcal{M}_{d,W}$} if for every $y\in\mathcal{M}_{d,W}$,
\begin{align*}
\Psi(\mathcal{W}_{W,K}(y))=\zero &\ \Longrightarrow\ y=y_{\mathrm{neu}},\\
\Psi(\mathcal{W}_{W,K}(y))=\nonzero &\ \Longrightarrow\ y\neq y_{\mathrm{neu}}.
\end{align*}
Its \emph{resolves set} is
\[
\operatorname{Res}(\Psi):=\{y\in\mathcal{M}_{d,W}:\Psi(\mathcal{W}_{W,K}(y))\neq \inconclusive\}.
\]
Given two sound procedures $\Psi_1,\Psi_2$ on $\mathcal{M}_{d,W}$, we say that $\Psi_1$ \emph{dominates} $\Psi_2$
if $\operatorname{Res}(\Psi_2)\subseteq \operatorname{Res}(\Psi_1)$ and
\[
\Psi_1(\mathcal{W}_{W,K}(y))=\Psi_2(\mathcal{W}_{W,K}(y))\qquad\text{for all }y\in \operatorname{Res}(\Psi_2).
\]
\end{definition}

\begin{definition}[Pipeline components]\label{def:pipeline-components}
Fix integers $W\ge 1$ and $K\ge 1$ and consider realizable window data $w\in\mathcal{W}_{W,K}(\mathcal{M}_{d,W})$.
We use the following (possibly partial) operations.
\begin{enumerate}[label=\textup{(\alph*)},leftmargin=2.2em]
\item \textbf{Aggregation / reconstruction ($A$).} On the identifiability locus from Theorem~\ref{thm:window-ident-global}, define
a partial map $A:\mathcal{W}_{W,K}(\mathcal{M}_{d,W})\dashrightarrow \mathcal{M}_{d,W}$ by letting $A(w)$ be the (locally) unique
signal $y$ (or parameter $\pi$) with $\mathcal{W}_{W,K}(y)=w$, when such a unique realization exists; otherwise $A(w)$ is undefined.
\item \textbf{Projection ($P$).} For a positive vector/configuration $x\in(\Rp)^n$ with $y=\log x$, set $P(y)$ as in
\eqref{eq:Pdef}. In the pipeline we apply this projection in log-coordinates to the reconstructed object.
\item \textbf{Coercive decision layer ($B$).} For a projected log-vector $u\in\R^n$ (typically $u=P(y)$), define the
certificate value
\[
\mathcal{B}(u):=\sum_{i=1}^n J\!\left(e^{u_i}\right),
\]
and define $B(u)\in\{\zero,\nonzero\}$ by $B(u)=\zero$ if $\mathcal{B}(u)=0$ and $B(u)=\nonzero$ otherwise.
\end{enumerate}
The associated finite-data decision map is evaluated as follows: given window data $w$, if $A(w)$ is defined then set $x:=A(w)$, form the projected log-vector $u:=P(\log x)$, and output $B(u)$; if $A(w)$ is undefined we return $\inconclusive$.

\begin{remark}[Forced $(P,B)$-dependence in the decision layer]\label{rem:forced-PB}
The ordering $A\to P\to B$ is not an arbitrary implementation choice. The defect notion and the zero-defect class are scale-invariant in log-coordinates: if $y=\log x$ then adding a constant multiple of $\mathbf 1$ changes $y$ but leaves $P(y)$ unchanged. Accordingly, any sound rule whose conclusions depend only on defect-type information can, without loss, be written as a map of the projected log-vector $u=P(\log x)$ (cf.\ Definition~\ref{def:defect} and Lemma~\ref{lem:Pkernel}). In our framework the coercive certificate value $\mathcal B(u)$ and the induced decision $B(u)\in\{\zero,\nonzero\}$ provide the canonical such decision core; thus any sound certification pipeline that reconstructs $x$ from window data may be expressed (up to notational equivalence) as first reconstructing $x$ (stage $A$), then projecting to $u=P(\log x)$ (stage $P$), and finally applying a defect-certifier (stage $B$).
\end{remark}

\begin{remark}[Non-redundancy of the stages]\label{rem:nonredundancy}
The three stages address distinct obstructions; in particular, none can be dropped without either (a) breaking the scale-invariance built into the defect notion, (b) losing quantitative robustness under noise, or (c) re-solving the inverse problem from window data.
\begin{enumerate}[label=(\alph*)]
\item (Need for $P$.) For any $c>0$ one has $P(\log(cx))=P(\log x)$, hence $\mathrm{Def}(cx)=\mathrm{Def}(x)$ by Definition~\ref{def:defect}. Any sound defect-based decision must therefore be invariant under global scaling. An implementation that feeds $\log x$ directly into the decision layer can violate this invariance, so centering via $P$ (or an equivalent projection) is essential.
\item (Need for $B$.) Even on the identifiability locus, noise perturbs $u$ by two-sided errors; the certification conclusions in Theorems~\ref{thm:eps-tolerant}--\ref{thm:eps-cert} rely on a coercive lower bound linking $\mathrm{Def}(x)$ to a computable quantity of $u$. Without such a bound one cannot guarantee membership in $\Mean_{2\eps}$ under perturbations.
\item (Need for $A$.) The decision layer is stated in terms of $u=P(\log x)$. To apply it to observed window data $w$, one must compute some invariant $u=\mathcal U(w)$ consistent with the model. On $\mathcal U_{d,W}$ this is achieved by $u=P(\log A(w))$ using the reconstruction map $A$ (Theorems~\ref{thm:window-ident-general} and \ref{thm:window-ident-global}). Any alternative sound pipeline must still implicitly recover the same invariant $u$ from $w$, i.e. it cannot bypass the inversion content of $A$.
\end{enumerate}
\end{remark}

For brevity we continue to denote this overall procedure by $\Phi^\ast=A\circ B\circ P$, with the understanding that $A$ is the reconstruction stage and that $\inconclusive$ is returned when reconstruction is not defined.
\end{definition}

\begin{proposition}\label{prop:pipeline-sound}
Fix $d\ge 0$, a window length $W\ge 1$, and $K\ge 2d+1$. Let $\mathcal M_{d,W}$ denote the degree-$\le d$
rational (finite-state) model class and let $\mathcal W_{W,K}:\mathcal M_{d,W}\to\R^K$ send a signal to its
first $K$ $W$-block window sums. On the nondegenerate identifiability locus singled out in
Theorem~\ref{thm:window-ident-global} (in particular, where the reconstruction step $A$ is well-defined by local invertibility / Hankel nondegeneracy), the canonical procedure
\[
\Phi^\ast \:=\ A\circ B\circ P
\]
is well-defined from $w=\mathcal W_{W,K}(y)$ and is sound in the sense of Definition~\ref{def:domination}.
\end{proposition}

\begin{proof}
The projection step $P$ is the mean-zero projection in log-coordinates (Definition~\ref{def:proj}).
By Definition~\ref{def:defect}, $\mathrm{Def}(x)=\|P(\log x)\|_2$, hence $\mathrm{Def}(x)=0$ if and only if $P(\log x)=0$.
Since $P$ is a projection, $P(P(\log x))=P(\log x)$. The coercive bound (Theorem~\ref{thm:coercive})
implies that vanishing projected cost forces $P(\log x)=0$, hence $\mathrm{Def}(x)=0$, giving soundness of the $B$ step.
Finally, on the identifiability locus, the aggregation step $A$ (reconstruction from $W$-block window sums)
does not introduce spurious neutral realizations: if the reconstructed projected cost is $0$, then the unique
signal compatible with the observed $w$ is neutral. Hence $\Phi^\ast$ is a sound finite-data procedure.
\end{proof}

\begin{remark}\label{rem:domination-intuition}
Soundness prevents false positives/negatives; domination compares how often a sound procedure can decide.
A dominating procedure is therefore ``best possible'' among sound procedures for the same finite data.
\end{remark}


\begin{lemma}\label{lem:sound-factors}

Let $\mathcal{U}_{d,W}\subset\mathcal{M}_{d,W}$ be a neighborhood of $y_{\mathrm{neu}}$ such that the restriction
$\mathcal{W}_{W,K}\vert_{\mathcal{U}_{d,W}}$ is injective (equivalently: for each realizable window vector
$w\in\mathcal{W}_{W,K}(\mathcal{U}_{d,W})$ there is at most one $y\in\mathcal{U}_{d,W}$ with $\mathcal{W}_{W,K}(y)=w$).
Let $\Psi:\R^K\to\{\zero,\nonzero,\inconclusive\}$ be sound on $\mathcal{M}_{d,W}$.
Then whenever $w$ is realized by some $y\in\mathcal{U}_{d,W}$, the value $\Psi(w)$ is forced solely by the
truth-value of the underlying predicate on that unique realization (neutral vs.\ non-neutral), and cannot
disagree with any other sound procedure on the same model class when restricted to $\mathcal{U}_{d,W}$.

\end{lemma}

\begin{proof}
If $w=\mathcal{W}_{W,K}(y)$ and $\Psi(w)=\zero$, soundness forces $y$ to be neutral; if $\Psi(w)=\nonzero$,
soundness forces $y$ to be non-neutral. Uniqueness of the realization for $w$ makes these implications
well-defined on the image of $\mathcal{W}_{W,K}$ and prevents contradictory sound outputs.
\end{proof}

We now state the \emph{local optimality} (domination) property of the canonical finite-data decision rule.
Write $\operatorname{Res}(\Psi)$ for the set of realizable instances on which a rule $\Psi$ returns a conclusive value in $\{\zero,\nonzero\}$.
We say that $\Phi$ \emph{dominates} $\Psi$ on a locus $\mathcal{V}$ if $\operatorname{Res}(\Psi)\cap\mathcal{V}\subseteq\operatorname{Res}(\Phi)\cap\mathcal{V}$ and the two rules agree on $\operatorname{Res}(\Psi)\cap\mathcal{V}$.
Informally, the rule first projects to the mean-zero subspace, then performs a window-based test, and finally outputs a ternary decision. The conclusion below is local: domination is asserted only on the identifiability neighborhood $\mathcal{U}_{d,W}$.

\begin{theorem}\label{thm:domination}
Fix integers $d\ge 0$, $W\ge 1$, and $K\ge 2d+1$.
Let $\mathcal{M}_{d,W}$ be a nondegenerate (open dense) parameter locus, and let
$\mathcal{U}_{d,W}\subset\mathcal{M}_{d,W}$ be a neighborhood of $y_{\mathrm{neu}}$ on which the restricted window map
$\mathcal{W}_{W,K}\vert_{\mathcal{U}_{d,W}}$ is injective.
Assume:


\begin{enumerate}[label=\textup{(\alph*)},leftmargin=2.2em]
\item for every $y\in\mathcal{U}_{d,W}$ the window-sum sequence $S_k:=\mathcal{W}(y)_k$ satisfies the hypotheses of Theorem~\ref{thm:window-ident-global}, so the data $(S_0,\dots,S_{2d-1})$ determine the induced window-sum process uniquely on $\mathcal{U}_{d,W}$; and
\item the realizable datum $w=\mathcal{W}_{W,K}(y)$ determines the neutral/non-neutral truth value on $\mathcal{U}_{d,W}$, i.e.
\[
\mathcal{W}_{W,K}(y)=\mathcal{W}_{W,K}(y_{\mathrm{neu}})\ \Longrightarrow\ y=y_{\mathrm{neu}}.
\]
\end{enumerate}

Let $\Phi^\ast:\mathcal{W}_{W,K}(\mathcal{M}_{d,W})\to\{\zero,\nonzero,\inconclusive\}$ denote the canonical procedure $A\circ B\circ P$ defined in this section.

Then for every sound procedure $\Psi:\R^K\to\{\zero,\nonzero,\inconclusive\}$ one has
$\operatorname{Res}(\Psi)\cap\mathcal{U}_{d,W}\subseteq \operatorname{Res}(\Phi^\ast)\cap\mathcal{U}_{d,W}$, and for every
$y\in\operatorname{Res}(\Psi)\cap\mathcal{U}_{d,W}$ the outputs agree on the realized datum $w=\mathcal{W}_{W,K}(y)$:
\[
\Psi(\mathcal{W}_{W,K}(y))=\Phi^\ast(\mathcal{W}_{W,K}(y)).
\]

\end{theorem}

\begin{proof}
Let $\Psi$ be any sound procedure on $\mathcal{M}_{d,W}$ and let $y\in\operatorname{Res}(\Psi)\cap\mathcal{U}_{d,W}$.
Write $w:=\mathcal{W}_{W,K}(y)$.


By assumption (a) and Theorem~\ref{thm:window-ident-global}, the realized datum $w$ determines the induced window-sum process (and hence its Prony parameters) uniquely on $\mathcal{U}_{d,W}$.
Assumption (b) guarantees that this information determines whether or not $y=y_{\mathrm{neu}}$ (equivalently: a neutral realization is compatible with $w$ if and only if $y=y_{\mathrm{neu}}$).

By construction, $\Phi^\ast$ evaluates exactly this truth value: it outputs $\zero$ exactly when the unique
reconstruction compatible with $w$ lies in the neutral \emph{zero-defect} class, and outputs $\nonzero$ exactly when
no neutral realization is compatible with $w$; otherwise it outputs $\inconclusive$.
Hence $\Phi^\ast$ is sound.

Now assume $y\in\operatorname{Res}(\Psi)\cap\mathcal{U}_{d,W}$. If $\Psi(w)=\zero$, then soundness forces $y=y_{\mathrm{neu}}$, so the unique
reconstruction from $w$ is neutral and therefore $\Phi^\ast(w)=\zero$.
If $\Psi(w)=\nonzero$, then $y\neq y_{\mathrm{neu}}$, so the unique reconstruction from $w$ is non-neutral and
$\Phi^\ast(w)=\nonzero$.
Thus whenever $\Psi$ resolves (in the sense of Definition~\ref{def:domination}), $\Phi^\ast$ resolves with the same answer, i.e.\ $\Phi^\ast$ dominates $\Psi$.
\end{proof}


\begin{corollary}[Global domination under global injectivity]\label{cor:domination-global}
Fix integers $d\ge 0$, $W\ge 1$, and $K\ge 2d+1$.
Assume:
\begin{enumerate}[label=\textup{(\alph*)},leftmargin=2.2em]
\item the window map $\mathcal{W}_{W,K}$ is injective on the parameter locus $\mathcal M_{d,W}$;
\item for every $y\in\mathcal M_{d,W}$, the induced window-sum sequence satisfies the hypotheses of Theorem~\ref{thm:window-ident-global};
\item the realizable datum determines global neutrality on $\mathcal M_{d,W}$, i.e.
\[
\mathcal{W}_{W,K}(y)=\mathcal{W}_{W,K}(y_{\mathrm{neu}})\ \Longrightarrow\ y=y_{\mathrm{neu}}
\quad\text{for all }y\in\mathcal M_{d,W}.
\]
\end{enumerate}
Then for every sound procedure $\Psi:\R^K\to\{\zero,\nonzero,\inconclusive\}$,
\[
\operatorname{Res}(\Psi)\subseteq \operatorname{Res}(\Phi^\ast),
\]
and for all $y\in\operatorname{Res}(\Psi)$ one has
\[
\Psi(\mathcal W_{W,K}(y))=\Phi^\ast(\mathcal W_{W,K}(y)).
\]
\end{corollary}

\begin{proof}
Apply Theorem~\ref{thm:domination} with $\mathcal U_{d,W}:=\mathcal M_{d,W}$.
\end{proof}



Fix a costed space $(S,O,c)$ with meaning map $\Mean(s)=\arg\min_{o\in O}c(s,o)$.
In applications one typically does not compute $\Mean(s)$ by scanning all $o\in O$.
Instead, one seeks a \emph{certificate functional} $\mathcal{C}$ that:

\begin{itemize}
\item can be evaluated efficiently from $(s,o)$ (often through a low-dimensional
 proxy or windowed measurements), and
\item is \emph{sound}: $\mathcal{C}(s,o)$ attains its minimum (or maximum) exactly
 when $o\in\Mean(s)$.
\end{itemize}

The simplest sound certificates are \emph{separable} in the sense that they depend
on $(s,o)$ only through a canonical ratio (or a monotone transform thereof), so
that the certificate is equivalent to comparing values of $J$.

\subsection{Canonical rescalings}\label{subsec:canon-rescale}

We make precise what is meant by a ``canonical rescaling'' of a signal in the present model.
Assume that the multiplicative group $(0,\infty)$ acts on $S$ (written $\lambda\cdot s$) and that
the scale map $\iota_S:S\to(0,\infty)$ is \emph{equivariant} for this action:
\[
\iota_S(\lambda\cdot s)=\lambda\,\iota_S(s)\qquad(\lambda>0,\ s\in S).
\]
We call $s' \in S$ a \emph{canonical rescaling} of $s\in S$ if $s'=\lambda\cdot s$ for some $\lambda>0$.
This captures the idea that $s$ and $s'$ represent the same ``configuration'' up to overall magnitude.

In several arguments below we also use that the penalty $J$ is strictly monotone on one branch.
Accordingly, for a fixed $s\in S$ we say that the candidate class $O$ lies on a \emph{single $J$-branch for $s$} if either
$\iota_S(s)/\iota_O(o)\ge 1$ for all $o\in O$ or $\iota_S(s)/\iota_O(o)\le 1$ for all $o\in O$.
On such a branch, $J$ is strictly monotone and hence invertible.

\subsection{Forced separability for canonical reciprocal costs}

We now state the main rigidity claim in the simplest form that is directly used
later. The proof is a short functional argument: the coercivity estimate together with the reciprocity structure induced by the cost assumptions pin down the allowable comparisons between candidates $o\in O$.
Such order-preserving representation arguments are classical in the theory of functional equations; see \cite{aczel1966}.


\noindent This is a standard order-theoretic fact; we include it for completeness.
\begin{lemma}\label{lem:forced-factor}
Let $c:S\times O\to[0,\infty)$ be the ratio-induced cost $c(s,o):=J(\iota_S(s)/\iota_O(o))$.
Fix a certificate $\mathcal{C}:S\times O\to\R$.
Assume that $\mathcal{C}$ is \emph{ordinally sound} relative to $c$ in the sense of
Definition~\ref{def:sound}(ii).
Then for each $s\in S$ there exists a strictly increasing function
$\phi_s:\operatorname{im}(c_s)\to\R$ (where $c_s(o):=c(s,o)$) such that
\[
\mathcal{C}(s,o)=\phi_s(c(s,o))\qquad\text{for all }o\in O.
\]
In particular, on each fixed $s$, $\mathcal{C}(s,\cdot)$ is a strictly monotone reparametrization of $c(s,\cdot)$.
\end{lemma}

\begin{proof}
Fix $s\in S$. Ordinal soundness means $\mathcal{C}_s$ and $c_s$ induce the same preorder on $O$.
Define $\phi_s$ on $\operatorname{im}(c_s)$ by $\phi_s(r):=\mathcal{C}_s(o)$ for any $o$ with $c_s(o)=r$; this is well-defined and strictly increasing, and yields $\mathcal{C}_s=\phi_s\circ c_s$.
\end{proof}


\begin{proposition}\label{prop:rescale-stability}
Assume the hypotheses of Lemma~\ref{lem:forced-factor}.
Fix $s\in S$ and assume $O$ lies on a single $J$-branch for $s$ and for every canonical rescaling $s'=\lambda\cdot s$
(in the sense of Section~\ref{subsec:canon-rescale}).
Then for each such $s'=\lambda\cdot s$ there exists a strictly monotone function $\psi_{s,s'}:\R\to\R$ such that
\[
\mathcal{C}(s',o)=\psi_{s,s'}(\mathcal{C}(s,o))\qquad\text{for all }o\in O.
\]
In particular, canonical rescalings preserve the certificate ordering up to a strictly monotone reparametrization on any single $J$-branch.
\end{proposition}

\begin{proof}
Let $s'=\lambda\cdot s$ with $\lambda>0$. Write $r(o):=\iota_S(s)/\iota_O(o)$ and $r'(o):=\iota_S(s')/\iota_O(o)=\lambda r(o)$.
By Lemma~\ref{lem:forced-factor} there exist strictly monotone functions $\phi_s,\phi_{s'}$ such that
$\mathcal{C}(s,o)=\phi_s(J(r(o)))$ and $\mathcal{C}(s',o)=\phi_{s'}(J(r'(o)))$ for all $o\in O$.

By the single-branch hypothesis, the restriction of $J$ to the relevant range of ratios is strictly monotone and hence invertible.
Let $J^{-1}$ denote this branch inverse.
Define a function $F_\lambda:[0,\infty)\to[0,\infty)$ by
\[
F_\lambda(t):=J\!\bigl(\lambda\,J^{-1}(t)\bigr).
\]
Since $J$ and $J^{-1}$ are strictly monotone on the branch, $F_\lambda$ is strictly monotone. Moreover,
for every $o\in O$ we have $J(r'(o))=J(\lambda r(o))=F_\lambda(J(r(o)))$.

Since $\phi_s:\operatorname{im}(c_s)\to \operatorname{im}(\mathcal C_s)$ is strictly monotone, it is a bijection onto its image, so the inverse $\phi_s^{-1}$ is well-defined on $\operatorname{im}(\mathcal C_s)=\{\mathcal C(s,o):o\in O\}$. In particular, $\phi_s(J(r(o)))=\mathcal C(s,o)$ lies in this domain for every $o\in O$.
Therefore
\[
\mathcal{C}(s',o)=\phi_{s'}(J(r'(o)))=\phi_{s'}(F_\lambda(J(r(o))))
=\underbrace{\bigl(\phi_{s'}\circ F_\lambda\circ \phi_s^{-1}\bigr)}_{=:~\psi_{s,s'}}\bigl(\phi_s(J(r(o)))\bigr)
=\psi_{s,s'}(\mathcal{C}(s,o)).
\]
The composition $\psi_{s,s'}=\phi_{s'}\circ F_\lambda\circ \phi_s^{-1}$ is strictly monotone, as required.
\end{proof}


\begin{proposition}[Primitive ratio hypotheses imply the uniqueness bundle]\label{prop:primitive-to-H}
Let $c(s,o):=J(\iota_S(s)/\iota_O(o))$ and $\mathcal C:S\times O\to\R$.
Assume:
\begin{enumerate}[label=\textup{(P\arabic*)},leftmargin=2.2em]
\item \textbf{Ratio dependence}: if
$\iota_S(s_1)/\iota_O(o_1)=\iota_S(s_2)/\iota_O(o_2)$, then
$\mathcal C(s_1,o_1)=\mathcal C(s_2,o_2)$.
\item \textbf{Cost-to-ratio determinacy on the realized set}: if
$c(s_1,o_1)=c(s_2,o_2)$, then
$\iota_S(s_1)/\iota_O(o_1)=\iota_S(s_2)/\iota_O(o_2)$.
\item \textbf{State-ratio collapse at fixed observation}: for all $s_1,s_2\in S$ and $o\in O$,
\[
\frac{\iota_S(s_1)}{\iota_O(o)}=\frac{\iota_S(s_2)}{\iota_O(o)}.
\]
\end{enumerate}
Then the hypotheses (H1) and (H2) of Theorem~\ref{thm:forced-factor-unique} hold.
\end{proposition}

\begin{proof}
(H1) follows from (P2) then (P1): equal costs imply equal ratios, hence equal certificate values.
For (H2), fix $s_1,s_2,o$.
By (P3), the ratios for $(s_1,o)$ and $(s_2,o)$ are equal; applying (P1) yields
$\mathcal C(s_1,o)=\mathcal C(s_2,o)$.
\end{proof}


\begin{remark}[On the strength of (P3)]
Since $\iota_O(o)>0$, hypothesis (P3) implies $\iota_S(s_1)=\iota_S(s_2)$ for all $s_1,s_2\in S$, i.e.\ $\iota_S$ is constant.
In this regime the ratio $r(o)=\iota_S(s)/\iota_O(o)$ and hence the cost $c(s,o)=J(r(o))$ are \emph{independent of $s$}.
Accordingly, (H2) is automatic and the ``cross-state'' content of Theorem~\ref{thm:forced-factor-unique} becomes immediate.
If a nontrivial bridge from primitives to (H1)--(H2) is intended, (P3) should be weakened.
\end{remark}





\begin{proposition}[Reducing \textup{(H1)} to ordinal soundness plus cross-state calibration]\label{prop:reduce-H1}
Assume that for each $s\in S$ the candidate class $O$ lies on a single $J$-branch for $s$
(as in Section~\ref{subsec:canon-rescale}), so that $o\mapsto c(s,o)$ is strictly monotone in the ratio
$\iota_S(s)/\iota_O(o)$.
If $\mathcal C$ is ordinally sound (Definition~\ref{def:sound}), then for every $s\in S$ there exists a strictly
increasing map $f_s:c_s(O)\to\R$ such that
\[
\mathcal C(s,o)=f_s\bigl(c(s,o)\bigr)\qquad(o\in O),
\]
where $c_s(O):=\{c(s,o):o\in O\}$.
Moreover, hypothesis \textup{(H1)} of Theorem~\ref{thm:forced-factor-unique} holds if and only if the family
$\{f_s\}_{s\in S}$ is \emph{state-independent on overlaps} in the sense that for all $s_1,s_2\in S$,
\[
f_{s_1}(r)=f_{s_2}(r)\qquad\text{for all }r\in c_{s_1}(O)\cap c_{s_2}(O).
\]
In particular, a sufficient condition for \textup{(H1)} is that $c_s(O)$ contains a common interval
$I\subset\R$ for all $s$ and $f_s|_I$ is independent of $s$.
\end{proposition}

\begin{proof}
Fix $s$. Ordinal soundness says that $\mathcal C_s$ and $c_s$ induce the same preorder on $O$.
Since $c_s$ is a real-valued function, this is equivalent to the existence of a strictly increasing
reparameterization $f_s$ on the range $c_s(O)$ with $\mathcal C_s=f_s\circ c_s$.
(The construction is canonical: define $f_s(r):=\mathcal C(s,o)$ for any $o$ with $c(s,o)=r$; this is well-defined
because $c(s,o_1)=c(s,o_2)$ implies $\mathcal C(s,o_1)=\mathcal C(s,o_2)$ by the ``$\Leftrightarrow$'' in ordinal soundness.)
For the characterization of \textup{(H1)}, note that if $c(s_1,o_1)=c(s_2,o_2)=r$ then
\[
\mathcal C(s_1,o_1)=f_{s_1}(r),\qquad \mathcal C(s_2,o_2)=f_{s_2}(r),
\]
so $\mathcal C(s_1,o_1)=\mathcal C(s_2,o_2)$ for all such pairs if and only if $f_{s_1}(r)=f_{s_2}(r)$ on every
overlap value $r\in c_{s_1}(O)\cap c_{s_2}(O)$.
The final sufficient condition is immediate.
\end{proof}

\begin{remark}[On \textup{(H2)} and rescaling]
Hypothesis \textup{(H2)} of Theorem~\ref{thm:forced-factor-unique} is a genuine \emph{state-rigidity} condition:
it asserts that the certificate is constant across $S$ at each fixed observation $o$.
In applications this typically follows from an explicit invariance or normalization principle for the certificate,
rather than from ordinal soundness alone.
\end{remark}


\begin{theorem}[Global forced factorization with uniqueness]\label{thm:forced-factor-unique}
Let $c:S\times O\to[0,\infty)$ be the ratio-induced cost and let $\mathcal C:S\times O\to\R$ be a certificate.
Assume:
\begin{enumerate}[label=\textup{(H\arabic*)},leftmargin=2.2em]
\item \textbf{Cross-state cost dependence}: (A convenient sufficient route to verify \textup{(H1)} from ordinal soundness plus a single-branch calibration is given in Proposition~\ref{prop:reduce-H1}.)
for all $s_1,s_2\in S$ and $o_1,o_2\in O$,
\[
c(s_1,o_1)=c(s_2,o_2)\ \Longrightarrow\ \mathcal C(s_1,o_1)=\mathcal C(s_2,o_2).
\]
\item \textbf{State rigidity}: for all $s_1,s_2\in S$ and $o\in O$,
\[
\mathcal C(s_1,o)=\mathcal C(s_2,o).
\]
\end{enumerate}
Fix $s_\ast\in S$ and define
\[
\mathcal I_c:=\{r\in\R:\exists s\in S,\exists o\in O,\ r=c(s,o)\}.
\]
Then:
\begin{enumerate}[label=\textup{(\alph*)},leftmargin=2.2em]
\item there exists a unique map $\phi:\mathcal I_c\to\R$ such that
\[
\mathcal C(s,o)=\phi(c(s,o))\qquad\text{for all }s\in S,\ o\in O;
\]
\item there exists a unique map $\psi:O\to\R$ such that
\[
\mathcal C(s,o)=\psi(o)\qquad\text{for all }s\in S,\ o\in O.
\]
\end{enumerate}
\end{theorem}

\begin{proof}
For (a), define $\phi$ on $\mathcal I_c$ by choosing, for each $r\in\mathcal I_c$, any pair $(s_r,o_r)$ with
$c(s_r,o_r)=r$ and setting $\phi(r):=\mathcal C(s_r,o_r)$.
Assumption (H1) makes this well-defined.
If $r=c(s,o)$, then (H1) gives
$\mathcal C(s,o)=\mathcal C(s_r,o_r)=\phi(r)$, so the representation holds.
If $\widetilde\phi$ is another map with the same representation property, then for any
$r\in\mathcal I_c$ and any witness $(s,o)$ with $c(s,o)=r$,
\[
\widetilde\phi(r)=\mathcal C(s,o)=\phi(r),
\]
hence $\widetilde\phi=\phi$.

For (b), define $\psi(o):=\mathcal C(s_\ast,o)$.
By (H2), for every $s\in S$ one has
$\mathcal C(s,o)=\mathcal C(s_\ast,o)=\psi(o)$.
If $\widetilde\psi$ also satisfies $\mathcal C(s,o)=\widetilde\psi(o)$ for all $s,o$, then
evaluating at $s=s_\ast$ gives
$\widetilde\psi(o)=\mathcal C(s_\ast,o)=\psi(o)$ for every $o$, so $\widetilde\psi=\psi$.
\end{proof}

\begin{corollary}[Forced-factorization uniqueness from primitive hypotheses]\label{cor:forced-factor-unique-primitive}
Under the assumptions of Proposition~\ref{prop:primitive-to-H},
the uniqueness conclusions (a) and (b) of Theorem~\ref{thm:forced-factor-unique} hold.
\end{corollary}

\begin{proof}
By Proposition~\ref{prop:primitive-to-H}, assumptions (H1)--(H2) in
Theorem~\ref{thm:forced-factor-unique} are satisfied, so the result follows directly.
\end{proof}

\begin{remark}[Standalone status of the uniqueness theorem]
Theorem~\ref{thm:forced-factor-unique} and
Corollary~\ref{cor:forced-factor-unique-primitive}
are proved entirely by the arguments in this section.
\end{remark}


\subsection{\texorpdfstring{$\eps$-optimal certification from windowed measurements}{Epsilon-optimal certification from windowed measurements}}\label{subsec:eps-opt}

We record an $\eps$-robust variant that is useful when only approximate ratios are
available (e.g., from windowed measurements) and the optimizer may not be unique.

\begin{definition}\label{def:eps-meaning}
For $\eps\ge 0$ define the $\eps$-meaning set by
\[
\Mean_\eps(s)
:=\left\{o\in O:\; c(s,o)\le \inf_{o'\in O}c(s,o')+\eps\right\}.
\]
\end{definition}

\begin{lemma}\label{lem:J-Lipschitz}
Let $0<a\le b<\infty$ and let $J(x)=\tfrac12(x+x^{-1})-1$ on $(0,\infty)$.
Then $J$ is Lipschitz on $[a,b]$ with constant
\[
L(a):=\frac12\Bigl(1+a^{-2}\Bigr),
\qquad\text{i.e.}\qquad
|J(x)-J(y)|\le L(a)\,|x-y| \ \ \text{for all }x,y\in[a,b].
\]
In particular, if $|\widehat r/r-1|\le \delta$ and $r\in[a,b]$, then
\[
|J(\widehat r)-J(r)|\le L(a)\,\delta\,b.
\]
\end{lemma}

\begin{proof}
For $x>0$ we have $J'(x)=\tfrac12(1-x^{-2})$, hence for $x\in[a,b]$,
\[
|J'(x)|\le \tfrac12\bigl(1+a^{-2}\bigr)=L(a).
\]
By the mean value theorem, $|J(x)-J(y)|\le \sup_{[a,b]}|J'|\,|x-y|\le L(a)|x-y|$.
If $|\widehat r/r-1|\le\delta$ and $r\in[a,b]$, then $|\widehat r-r|\le \delta r\le \delta b$, giving the final bound.
\end{proof}
 % end lemma
\begin{theorem}\label{thm:eps-cert}

Assume the hypotheses of Lemma~\ref{lem:forced-factor}. Fix $s\in S$ and suppose that the relevant ratios
\[
r(o):=\iota_S(s)/\iota_O(o)
\]
lie in a known band $[a,b]\subset(0,\infty)$ for all candidates $o$ under consideration. Suppose we observe $K$
windowed measurements that allow us to estimate $r(o)$ by $\widehat r(o)$ with relative error at most $\delta$,
uniformly in $o$, i.e.\ $|\widehat r(o)/r(o)-1|\le \delta$. Then the proxy ranking obtained by replacing $r(o)$ with $\widehat r(o)$, i.e. by minimizing the perturbed costs $\widehat c(o):=J(\widehat r(o))$, can only localize candidates within $\Mean_{2\eps}(s)$, where one may take the explicit bound
\[
\eps \;=\; L(a)\,\delta\, b,
\qquad L(a):=\tfrac12(1+a^{-2}),
\]
so in particular $\eps(\delta)\to 0$ as $\delta\to 0$ (for fixed $a,b$).

\end{theorem}

\begin{proof}
Let $\widehat r(o)$ be the estimated ratio for candidate $o$, so that $|\widehat r(o)/r(o)-1|\le \delta$ where
$r(o)=\iota_S(s)/\iota_O(o)$. Since $r(o)\in[a,b]$ by hypothesis, Lemma~\ref{lem:J-Lipschitz} gives the uniform bound

\[
|J(\widehat r(o))-J(r(o))|\le \eps
\qquad\text{with}\qquad
\eps=L(a)\,\delta\,b.
\]

Thus ranking candidates using the perturbed costs $J(\widehat r(o))$ may change the true costs by at most $\eps$ in each direction. If $\widehat o$ minimizes the perturbed cost $\widehat c(o):=J(\widehat r(o))$, then for any $o\in O$,
\[
c(s,\widehat o)\le \widehat c(\widehat o)+\eps\le \widehat c(o)+\eps\le c(s,o)+2\eps.
\]
Hence any sound selection based on these measurements can only certify candidates within $\Mean_{2\eps}(s)$ (Definition~\ref{def:eps-meaning}).
\end{proof}


\begin{remark}[Standalone status of the $\eps$-noise bound]
Theorem~\ref{thm:eps-cert} follows directly from
Lemma~\ref{lem:J-Lipschitz} and the perturbation estimate proved above.
\end{remark}


\section{Applications}\label{sec:applications}

We illustrate the framework on simple model families and finite-window measurements. The examples show how the quantitative bounds translate into explicit tolerances and how the defect-based certificate behaves on neutral versus non-neutral data. We also note practical considerations for implementation and parameter choice.



\paragraph{A worked numerical instantiation (toy).} Take the scalar band $[a,b]=[1/2,2]$ so that $L(a)=\tfrac12(1+a^{-2})=\tfrac12(1+4)=\tfrac52$. If the relative ratio-estimation error satisfies $|\widehat r/r-1|\le \delta=0.01$, then Theorem~\ref{thm:eps-cert} yields $\eps=L(a)\,\delta\,b=(\tfrac52)(0.01)(2)=0.05$, and hence any minimizer of the perturbed cost is certified only up to the tolerance set $\Mean_{2\eps}(s)=\Mean_{0.1}(s)$. For example, for candidates with true ratios $r_1=1$ and $r_2=1.08$, we have $J(r_1)=0$ and $J(r_2)=\tfrac12(1.08+1/1.08)-1\approx 0.00148$, so the defect gap is far below $0.1$ and cannot be resolved at this noise level.


Based on the main theorems, a finite-data certification map $\Phi^\ast = A\circ B\circ P$
supports short-window identification of ``zero-defect'' (neutral) configurations for positive
signals. The framework naturally applies in settings where observations are aggregated over
short windows and one seeks a reliable, explainable decision rule derived from the certified
structure. The following examples illustrate representative domains and typical deployment
goals.

\subsection{A toy worked example (synthetic)}
To illustrate the objects (without claiming a realistic data model), take $W=8$ and a rational degree-$1$ sequence
$y$ generated by $\theta(z)=\frac{1}{1-\alpha z}$ with $|\alpha|<1$, so $y_n=\alpha^n$. Then the $k$th window sum is a geometric-series computation:
\[
\mathcal{W}(y)_k
=\sum_{j=0}^{7} \alpha^{8k+j}
=\alpha^{8k}\sum_{j=0}^{7}\alpha^{j}
=\alpha^{8k}\,\frac{1-\alpha^{8}}{1-\alpha}\qquad (\alpha\neq 1).
\]
(If $\alpha=1$, then $y_n\equiv 1$ and $\mathcal W(y)_k=8$ for all $k$.)
\paragraph{Tracing $A\circ B\circ P$ in this toy case.}
Fix $\alpha=\tfrac12$. Then
\[
\mathcal W(y)_0=\sum_{j=0}^7 2^{-j}=\frac{1-2^{-8}}{1-2^{-1}}=\frac{255}{128},
\qquad
\mathcal W(y)_1=2^{-8}\mathcal W(y)_0=\frac{255}{32768}.
\]
Thus the reconstruction map $A$ is explicit: for any $k$ with $\mathcal W(y)_k\neq 0$ one has
$\alpha^8=\mathcal W(y)_{k+1}/\mathcal W(y)_k$, hence $\alpha=(\mathcal W(y)_{k+1}/\mathcal W(y)_k)^{1/8}$.
Given a noisy observation $\widehat S=(\widehat S_0,\widehat S_1,\dots)$, one forms log-ratios and applies
$P$ to remove the global multiplicative scale, then uses $B$ to score candidates by the canonical cost $J$.

For two candidate parameters $\alpha,\beta$ (hence two candidates $o_\alpha,o_\beta$) the ratio of predicted window sums
over a fixed block determines the ratio band $[a,b]$ used in Theorem~\ref{thm:eps-cert}. The defect gap
$\Delta(s,o)=c(s,o)-\inf_{o'}c(s,o')$ then certifies the closer candidate in the $J$-cost, and the $\eps$-meaning set
quantifies robustness when the observed window sums are perturbed.

\subsection{Streaming anomaly detection and quality control}
Many industrial or operational systems expose \emph{positive} telemetry (rates, counts, intensities)
whose nominal state is well-modeled as a neutral baseline up to multiplicative drift.
Examples include packet or transaction counts, sensor intensities, and throughput rates.
Given a stream $t\mapsto x(t)\in(\Rp)^n$ sampled at high rate, one may form windowed aggregates
in the $W$-window regime (e.g. $W=8$) treated in Section~\ref{sec:aggregation} and compute the certificate $\Phi^\ast$
on each window.
Soundness guarantees that ``certificate $=0$'' implies genuine neutrality for the modeled class,
while the coercivity step $B$ provides a quantitative lower bound converting measured cost into defect.
This yields a principled decision rule: flag a window when the certified defect exceeds a chosen tolerance.
In practice one can maintain an online baseline by applying the $J$-projection $P$ (mean subtraction in log-coordinates)
to remove multiplicative bias, then monitoring the aggregated coercivity bound.

\subsection{Sensor normalization under multiplicative gain and drift}
Calibration problems often separate into (i) removing multiplicative gain/drift and (ii) certifying that
the remaining configuration is ``balanced''.
If a sensor array outputs positive readings $y=g\odot x$ with unknown gain vector $g>0$ and latent state $x>0$,
then in log-coordinates the gain becomes additive.
The projection step $P$ is designed to remove the neutral (mean) component in log-coordinates,
so it can be used as a normalization primitive prior to downstream certification.
When a protocol collects repeated short-window measurements, the aggregation map $A$
permits certification from window sums without reconstructing full-resolution trajectories.
The resulting workflow is: apply $P$ to normalize, apply $B$ to convert residual cost into a defect bound,
and apply $A$ to assemble a window-level certificate that is stable under the stated nonexpansive property
in log-coordinates.

\subsection{Low-latency detection in communications and network monitoring}
In networked systems, one often observes positive quantities (link utilizations, loss rates, queue occupancies)
and must detect departures from a neutral operating regime with minimal delay.
The forced factorization $\Phi^\ast=A\circ B\circ P$ is well-suited to low-latency operation:
$P$ is a simple local transformation (mean subtraction in log-space),
$B$ is a closed-form coercivity bound with sharp constant,
and $A$ uses short windows (here $W$ is fixed, e.g. $W=8$) to produce a decision statistic.
When the signal class is well-approximated by a finite-state (rational/recurrence) model,
Remark~\ref{rem:generic-uniqueness} explains why one expects generic identifiability from $K=2d+1$ windows,
making the procedure practically ``complete'' outside degenerate parameter regimes where the stated full-rank hypothesis fails.
This provides an interpretable alternative to black-box detectors: the certificate is derived from axioms and
enjoys a provable optimality property among sound finite-data strategies.

\subsection{Model-based diagnostics in biomedical and neuroscience pipelines}
Many biomedical measurements are nonnegative (spike counts, fluorescence intensities, normalized power spectra). In applications with zeros or sign changes, log-coordinates are not directly available and one must either restrict to strictly positive regimes or introduce an offset/likelihood model before applying the defect geometry.
In settings where short-window summaries are the primary data product (e.g.\ binned spike counts),
the aggregation-based theory provides a direct bridge between measurements and neutrality certification.
A common diagnostic task is to detect whether a regime is ``balanced'' relative to a subject-specific baseline
while being robust to multiplicative scale differences.
The projection $P$ addresses scale/baseline, while coercivity $B$ yields a defect lower bound that can be used to
compare windows across subjects or sessions.
As above, all guarantees remain conditional on the modeling hypotheses (positivity, admissible window maps,
and the specified nondegeneracy regime).

\subsection{Reliability screening in inverse problems and compressed summaries}
In inverse problems one frequently works not with raw signals but with compressed or summarized observations
(rolling averages, block sums, local histograms), especially when storage or privacy constraints are present.
The window-aggregation results show that certification can be performed directly on such summaries
without reconstructing the underlying full-resolution positive vector, provided the summary map matches the
aggregation model.
This can be used as a \emph{reliability screen}: before attempting a costly reconstruction,
one first checks whether the summary data certify neutrality (defect zero) or quantify the minimal defect
consistent with the summary.
If the certificate indicates large defect, one can either reject the instance or route it to a different solver;
if it indicates neutrality, one gains a rigorous justification for treating the instance as defect-free in downstream analysis.




\appendix
% Single appendix: display unlettered heading, keep internal A.* numbering
\refstepcounter{section}\label{app:RCL}
\section*{Appendix}
\addcontentsline{toc}{section}{Appendix}
\setcounter{subsection}{0}
\subsection*{The Recognition Composition Law: axioms and uniqueness}

In this appendix we record the self-contained consequences of the Recognition Composition Law used in the paper. Under the stated axioms and regularity assumptions, we derive normalization and reciprocity, and then prove the corresponding uniqueness statement for the reciprocal cost function.

\subsection{Standing axioms and assumptions}\label{app:RCL:axioms}


\paragraph{Axiom 1 (Recognition Composition Law).} For all $x,y>0$,
\begin{equation}\tag{RCL}\label{eq:RCL_app}
J(xy)+J(x/y)=2J(x)+2J(y)+2J(x)J(y).
\end{equation}

\paragraph{Axiom 2 (Calibration in log-coordinates).} Define $J_{\log}(t):=J(e^{t})$. We assume $J_{\log}$ is twice differentiable at $t=0$ and satisfies $J_{\log}''(0)=1$.


\subsection{Derived normalization and reciprocity}

\begin{proposition}[Normalization is forced]\label{prop:norm-forced-vv}

If $J:(0,\infty)\to\R$ is non-constant and satisfies the RCL \eqref{eq:RCL_app}, then $J(1)=0$.

\end{proposition}

\begin{proof}

Setting $x=y=1$ in \eqref{eq:RCL_app} gives
\[
2J(1)=2J(1)^2+4J(1),
\]
so $J(1)^2+J(1)=0$ and therefore $J(1)\in\{0,-1\}$.
If $J(1)=-1$, substituting $y=1$ into \eqref{eq:RCL_app} yields $J(x)\equiv -1$ for all $x>0$, contradicting non-constancy.
Hence $J(1)=0$.

\end{proof}

\begin{corollary}[Reciprocity is forced]\label{cor:recip-forced-vv}

Under the same hypotheses, $J(x)=J(1/x)$ for all $x>0$.

\end{corollary}

\begin{proof}

Set $x=1$ in \eqref{eq:RCL_app} and use Proposition~\ref{prop:norm-forced-vv}.

\end{proof}

\subsection{Main uniqueness theorem}\label{app:RCL:uniqueness}

\begin{theorem}\label{thm:RCL-unique}

Assume that $J:(0,\infty)\to\mathbb{R}$ is continuous, non-constant, and satisfies Axioms~1--2 of Section~\ref{app:RCL:axioms}. Then the cost is forced to be the canonical reciprocal form

\[
J(x)=\tfrac12\bigl(x+x^{-1}\bigr)-1=\frac{(x-1)^2}{2x}.
\]

In particular, the closed form is not an additional modeling choice: it is determined by the composition identity together with the local quadratic calibration at balance.

\end{theorem}

\begin{proof}
Define $g(x):=1+J(x)$ and pass to log-coordinates by setting $h(t):=g(e^{t})$. A direct substitution of \eqref{eq:RCL_app} shows that $h$ satisfies the classical d'Alembert functional equation
\[
h(s+t)+h(s-t)=2h(s)h(t)\qquad (s,t\in\mathbb{R}).
\]
By Proposition~\ref{prop:norm-forced-vv} we have $J(1)=0$, hence $h(0)=g(1)=1$. The classification of continuous solutions to d'Alembert's equation with $h(0)=1$ is standard (see, for example, \cite{aczel1966}): there exists $\lambda\ge 0$ such that either $h(t)=\cos(\lambda t)$ or $h(t)=\cosh(\lambda t)$.

The calibration Axiom~2 says that $J_{\log}(t)=J(e^{t})$ has second derivative $1$ at $t=0$. Since $J_{\log}(t)=h(t)-1$, this gives $h''(0)=1$. For the cosine branch one has $h''(0)=-\lambda^2\le 0$, which is incompatible with $h''(0)=1$. Hence we are in the hyperbolic branch and $h(t)=\cosh(\lambda t)$. Imposing $h''(0)=\lambda^2=1$ yields $\lambda=1$, so $h(t)=\cosh(t)$.

Finally, returning to multiplicative coordinates, $g(x)=h(\ln x)=\cosh(\ln x)=\tfrac12(x+x^{-1})$ and therefore $J(x)=g(x)-1=\tfrac12(x+x^{-1})-1=\frac{(x-1)^2}{2x}$.
\end{proof}


\subsection{Derived properties}\label{app:RCL:derived}


\begin{lemma}[Calibration at $1$ vs.\ $0$ in log-coordinates]\label{lem:calibration-equivalence}
Assume $J$ is differentiable at $1$ and satisfies reciprocity $J(x)=J(1/x)$. Then $J'(1)=0$.
If moreover $J$ is twice differentiable at $1$, then for $J_{\log}(t):=J(e^t)$ one has
$J_{\log}''(0)=J''(1)$. In particular, the condition $J_{\log}''(0)=1$ is equivalent to $J''(1)=1$.
\end{lemma}

\begin{proof}
Differentiating $J(x)=J(1/x)$ at $x=1$ gives $J'(1)=-J'(1)$, hence $J'(1)=0$.
By the chain rule, $J_{\log}'(t)=J'(e^t)e^t$ and $J_{\log}''(t)=J''(e^t)e^{2t}+J'(e^t)e^t$.
Evaluating at $t=0$ and using $J'(1)=0$ yields $J_{\log}''(0)=J''(1)$.
\end{proof}



Once the closed form is established, standard structural properties become immediate: the cost is reciprocal ($J(x)=J(1/x)$), nonnegative with equality only at $x=1$, strictly convex in log-coordinates near balance, and it diverges as $x\to 0^{+}$ or $x\to+\infty$. The main paper uses only the explicit formula and the fact that it is forced by Axiom~1, non-constancy, and the local calibration in Axiom~2.


\section{\texorpdfstring{Finite-field Jacobian matrix for $(d,W)=(3,8)$}{Finite-field Jacobian matrix for (d,W)=(3,8)}}\label{app:jacobian-38}
For completeness, we record the integer matrix $DF_{3,8}(\pi_0)$ used in the finite-field determinant check in Lemma~\ref{lem:witness-38}.

\[
\resizebox{\textwidth}{!}{$
	DF_{3,8}(\pi_0)=
	\begin{pmatrix}
		1& 7& -3& -9& 69& -54& 3\\
		0& 1008& -956& -1388& 25920& -20844& 5568\\
		0& 175456& -190016& -200544& 5088848& -4927712& 1787216\\
		0& 28857344& -34167296& -27911296& 755009920& -959230336& 430318144\\
		0& 4538167296& -5754321920& -3726310400& 82750894080& -165775961088& 89416058880\\
		0& 686488772608& -922559823872& -473024225280& 3615042621440& -26067531849728& 16866417889280\\
		0& 100127404457984& -141965037535232& -56107543330816& -1341867460689920& -3737714520064000& 2956546669428736
	\end{pmatrix}
	$}
\]

\section*{Conclusion}
We developed a finite-data framework that certifies whether a positive configuration is neutral (zero defect) from a short list of aggregated observations. The main contribution is a canonical decision pipeline $\Phi^\ast=A\circ B\circ P$ with distinct roles: $P$ removes scale ambiguity by projecting to mean-zero log-coordinates; $B$ turns canonical reciprocal cost into an explicit defect bound via a sharp coercivity inequality; and $A$ reconstructs the induced window-sum process on a standard nondegenerate locus (local invertibility, equivalently Hankel/Prony nondegeneracy). On this locus, the first $K$ window sums uniquely determine the window process, yielding sound and locally maximal certification: any sound procedure that resolves must agree with $\Phi^\ast$ near the neutral realization. We also introduce an $\varepsilon$-tolerant noise model and stability bounds that quantify how small window perturbations affect certified defect guarantees.

These results clarify what finite windows can—and cannot—decide. Reconstruction requires nondegeneracy; outside that regime, the correct outcome is “inconclusive,” not a potentially misleading recovery claim. Future directions include extending reconstruction beyond separable exponential-sum models, strengthening robustness under weaker separation assumptions, and developing practical guidelines for monitoring when only pooled or privacy-preserving aggregates are available.



\section*{Acknowledgements}
The authors thank Philip Beltracchi, Sebastian Pardo-Guerra, and Milan Zlatanovic for careful reading and detailed comments and suggestions during the preparation of this manuscript. Their feedback improved the clarity, notation, and overall structure, and strengthened the quality of the paper. The authors received no external funding.

\begin{thebibliography}{10}

\bibitem{aczel1966}
J.~Acz\'el.
\newblock \emph{Lectures on Functional Equations and Their Applications}.
\newblock Academic Press, 1966.

\bibitem{washburnbarghi2026RCC}
Washburn, J.; Rahnamai Barghi, A.
\newblock Reciprocal Convex Costs for Ratio Matching: Axiomatic Characterization.
\newblock \emph{Axioms} \textbf{2026}, \emph{15}(2), 151. doi:10.3390/axioms15020151.

\bibitem{washburnnetal2026RG}
Washburn, J.; Zlatanovi\'c, M.; Allahyarov, E.\ \emph{Recognition Geometry}. \emph{Axioms} 2026, 15, 90.

\bibitem{moreau1965}
J.-J.~Moreau.
\newblock Proximit\'e et dualit\'e dans un espace hilbertien.
\newblock \emph{Bull. Soc. Math. France}, 93:273-299, 1965.

\bibitem{bregman1967}
L.~M.~Bregman.
\newblock The relaxation method of finding the common point of convex sets and its application to the solution of problems in convex programming.
\newblock \emph{USSR Comput. Math. Math. Phys.}, 7(3):200-217, 1967.

\bibitem{rockafellar1976}
R.~T.~Rockafellar.
\newblock Monotone operators and the proximal point algorithm.
\newblock \emph{SIAM J. Control Optim.}, 14(5):877-898, 1976.


\bibitem{osborne1995prony}
M.~R.~Osborne and G.~K.~Smyth.
\newblock A modified {P}rony algorithm for exponential function fitting.
\newblock \emph{SIAM J. Sci. Comput.}, 16(1):119--138, 1995. doi:10.1137/0916008.

\bibitem{peller2003hankel}
V.~V.~Peller.
\newblock \emph{Hankel Operators and Their Applications}.
\newblock Springer Monographs in Mathematics. Springer-Verlag, New York, 2003.

\end{thebibliography}
\end{document}
